\documentclass[11pt,a4paper]{article}
\PassOptionsToPackage{table}{xcolor}

% ============================================================
%  The Sweet Cycle  (Paper XII)
%  Philosophy of Intimacy and the Theory of Justice
%  Uses the shared series style (serendip-paper.sty). XeLaTeX +
%  Noto Serif CJK SC for Chinese terms. House rules: no em dash.
% ============================================================

\usepackage{fontspec}
\usepackage{serendip-paper}
\usepackage{ragged2e}

% Paper-specific colours used in the body, mapped into the series palette.
% (rose, gold, ink are provided by serendip-paper.)
\definecolor{deeprose}{HTML}{c56683}    % the series' rosedeep
\definecolor{blushground}{HTML}{fdf8f9} % the series' page (pageblush)

% Inline CJK helper (\cjkfont is provided by serendip-paper).
\newcommand{\cn}[1]{{\cjkfont #1}}

% Gold rule for section breaks.
\newcommand{\goldrule}{\medskip\noindent{\color{gold}\rule{\linewidth}{0.4pt}}\medskip}

\begin{document}

% ---------- Title page ----------
\thispagestyle{empty}
\null\vfill
\begin{center}
{\footnotesize\color{rosedeep}\scshape Philosophy of Intimacy and the Theory of Justice \quad$\cdot$\quad Paper XII\par}
\vspace{0.9em}
\coverrule\par
\vspace{1.6em}
{\color{warnred}\Large The Sweet Cycle\par}
\vspace{0.5em}
{\color{warnred}\cjkfont\large 甜之循环的认识与正义实践论\par}
\vspace{0.25em}
{\color{rose}\cjkfont\normalsize 神经科学、精神分析与政治经济学的讨论\par}
\vspace{2.4em}
\coverrule[0.25]\par
\vspace{1.4em}
{\large Wanhong Huang\par}
\vspace{0.5em}
{\footnotesize \texttt{huangwanhong@serendip.ngo}\par}
\vspace{1.2em}
{\today\par}
\end{center}
\vfill\null
\coverfootnote
\clearpage

% ============================================================
\begin{abstract}
\noindent
On a magazine cover in Ginza the author encountered a proposition stated without hedging: that sweetness is justice. The proposition is the occasion of this paper, not its target. What it names, in the idiom of intimate life, is a cycle that appears to harm no one and to benefit everyone who can be seen within it. A firm sells, a person who is indulged feels cared for, a person who indulges feels needed, the bond is strengthened, and the strengthened bond returns as a willingness to buy again. Inside its own boundary the cycle closes upon a surplus of satisfaction, and the satisfaction is real. This paper asks on what grounds, and under what conditions, such a cycle may be judged genuinely good rather than a forged closure that sustains its inner surplus by extraction across its outer edge. The inquiry proceeds by placing several incommensurable frameworks before the same phenomenon, psychoanalysis, neuroscience, the ethics of care, political economy, and feminist theory, and letting each say what it can and cannot say. None is permitted to adjudicate the others. A psychoanalytic foundation, reconstructed here in full from the author's prior work on lack, the \textit{objet a}, and the cute as a mediation of lack, supplies the structural diagram on which the later analysis is fought. The political economy of \textit{amai}, in which capital does not merely serve a pre-existing desire but produces and reproduces the desire itself and extracts a surplus from its reproduction, receives the most sustained treatment, because it is there that the difference between a genuine and a forged cycle becomes measurable. The paper does not close on a single verdict. It closes on a plurality of practices, just, ethical, eudaimonic, and co-creative, by which a sweet cycle might be made to keep its value within itself and to return that value to those who make it, with no one reduced to fuel.
\end{abstract}

\goldrule

% ============================================================
\section{Prologue: Sweetness as Justice in Ginza}

On a magazine cover in Ginza the author read a phrase that arranged cosmetics, small accessories, and clothing under a single heading: that the sweet is justice. \cn{甘いは正義。} I take this phrase as it presents itself in public, as a proposition, and I do not pass through it to the contents it advertises or to the people who make or read it. The proposition is enough. It is enough because it performs, in four characters, a movement that this paper will spend many pages trying to slow down and examine: it takes a register of intimate life, sweetness, indulgence, the small daily economy of being charmed and charming, and it confers on that register the name of a virtue. Not pleasure, not comfort, but justice. The naming is doing something, and the first task is to feel why it is at once so persuasive and so unsettling.

It is worth pausing on the grammar of the phrase before pausing on its truth, because the grammar is already an argument. The copula is doing the work. To say that the sweet is pleasant would be to report a feeling; to say that the sweet sells would be to report a fact about markets; to say that the sweet is kind would be to praise a disposition. The phrase says none of these. It places sweetness and justice on the two sides of an identity, and an identity claim of this form does something a mere predication does not: it offers to settle in advance a question that justice, of all concepts, exists precisely to keep open. Justice is the name we give to the unfinished work of asking whether an arrangement is defensible to all whom it touches. To declare that something simply \textit{is} justice is to propose that the asking is over, that the verdict is in, that one may now enjoy the arrangement with the conscience at rest. The cover, in other words, does not merely advertise objects. It advertises a kind of permission. And the discomfort that prompts this paper is, at bottom, a discomfort with permission granted too quickly, with a verdict announced before the trial.

This is the deeper reason the proposition cannot be passed over as a mere slogan. Slogans that promise pleasure or beauty make no claim on the conscience and so provoke no scruple. A slogan that promises justice enlists the conscience on the side of consumption, and it is the enlistment, not the consumption, that asks to be examined. The phrase is an unusually honest specimen precisely because it says the quiet part aloud. Much of consumer culture works by keeping its appeal in the register of desire and never raising the question of justice at all. This cover raises it, names it, and answers it in the same breath, and in doing so it hands the critic both the question and a confession that the question was felt to need answering.

It is persuasive because the cycle it names does seem, from inside, to be good. Consider the cycle in its plainest form. A person buys something that makes them feel sweet, cute, worth indulging. They bring this sweetness into a relationship, where it is received as an invitation to care. The one who cares feels needed, and being needed is its own reward. The bond tightens. A tighter bond makes both parties more willing, next time, to spend on the small mediations that let the sweetness continue. The firm that sold the first object profits and is thereby enabled to offer the next. Every party named in this description ends better off than they began. No one inside the circle is visibly wronged. If justice is, in one serviceable sense, an arrangement in which value is generated and returned to those who generate it, then the cycle appears to satisfy it. This is not a strawman to be knocked down. It is, I will argue at length, the strongest thing that can be said for \textit{amai}, and it must be said in full strength before anything else is said.

It will help to give the cycle a face, since abstraction tends to make these loops look thinner than they are in life. Imagine a person who, on a difficult morning, chooses a small bright thing, a tinted balm, a ribbon, a charm for a bag, and feels, in choosing it, a lift that is not vanity but something closer to self-address, a way of saying to oneself that one is worth a small tenderness. That tenderness is carried into the day and into a relationship, where the partner notices, is charmed, and answers the charm with an attention that would not have been called forth otherwise. The person feels met. The partner feels the particular warmth of being the one who meets. Neither would describe what passed between them as a transaction, and they would be right, because what passed between them was recognition, and recognition is not bought even when its occasion was. The next time the small bright thing is chosen, it is chosen against the memory of having been met, and so the object now carries a freight of relation it did not carry at first. This is not a trivial circuit. It is one of the ways ordinary love actually sustains itself across the wear of ordinary time, and any account that cannot feel its genuine warmth has disqualified itself from judging it.

One can put the cycle's apparent goodness in the formal vocabulary the series has developed, and doing so raises rather than lowers the stakes. The loop does not look like sterile repetition, the mere going-around that returns each party to where they began with nothing gained. It looks like a spiral that climbs: each turn deepens the bond a little beyond the last, something is laid down that was not there before, the return is not to the same point but to a higher one. In the series' terms the cycle appears to carry a positive holonomy, a real accumulation of relational value across each traversal rather than a flat circulation that nets to nothing. And it appears to carry that holonomy from within, generated by the participants' own exchange and not imported from any reservoir outside them. Were that appearance to hold under scrutiny, the cycle would be not merely permissible but exemplary, a small working proof that intimate life and economic life can be made to feed one another without either degrading the other. The whole drama of this paper lies in the words \textit{were that appearance to hold}.

It is unsettling for a reason that the cover itself supplies. The cover does not say that sweetness is pleasant, or profitable, or kind. It says that sweetness is just. And a claim of justice is never only a description of those who are present in the frame. Justice is a claim about a whole, about who is counted and who is not, about where the costs of an arrangement come to rest. The moment sweetness is named as justice, the question of justice is opened, and that question cannot be answered by looking only at the satisfied faces inside the circle. It requires us to ask what lies across the circle's edge, who or what is reproducing the conditions of all that sweetness, and whether the inner surplus is generated within the loop or drawn in from outside it. The cover, by reaching for the word justice, hands us the very standard by which its own proposition must be tested.

This paper is the twelfth in a series on the philosophy of intimacy and the theory of justice, and it occupies a particular place in that series. Where earlier papers built the formal vocabulary, the distinction between a cycle that accumulates value within itself and one that only appears to, the criterion of return to the creator, the diagnosis of an open system that masquerades as a closed loop, this paper does not build that vocabulary again. It puts it to work on a single, concrete, contemporary phenomenon. \textit{Amai} is not the goal of the inquiry. It is the inquiry's best instrument, because it is a case in which the reality of the satisfaction and the suspicion of the structure are both maximal at once. The relationship really is strengthened. The labor that underwrites the strengthening may really be hidden. A case where one of these were weak would teach us little. \textit{Amai} is sharp precisely because neither is weak.

Two terms from the series will recur often enough to deserve a plain statement here, so that a reader who has not followed the earlier papers is not left guessing. The first is holonomy, borrowed from the geometry of transport around a loop. The image is of carrying something, a value, an orientation, a measure of worth, once around a closed path and asking whether it returns changed. If it returns to its starting point exactly as it left, the loop has done no lasting work, and the series calls this a zero holonomy, the signature of mere circulation. If it returns enriched, with a real increment that the traversal itself produced, the loop has accumulated, and the series calls this a positive holonomy, the signature of genuine generation. The second term is the forged holonomy, and it names the case that this paper will argue \textit{amai} exemplifies: a loop that displays the increment of a positive holonomy but did not generate it internally, having instead drawn it in from outside the path while presenting itself as closed. The whole question of whether a sweet cycle is good reduces, in these terms, to whether its evident accumulation is generated or forged, and that is a question one cannot answer without following the value across the edge the cycle would prefer we not cross.

A word on method, since the method is itself a claim. I will not resolve \textit{amai} into a verdict by subordinating every framework to a master framework that pronounces the truth of the rest. I will instead set several frameworks before the same phenomenon, psychoanalysis, neuroscience, the ethics of care, political economy, feminist theory, and require each to state, in its own voice, what it is able to see and what it is constitutively unable to see. They will not agree. Psychoanalysis will describe a desire that no satisfaction closes; neuroscience will confirm that the satisfaction is a real event in the nervous system; the ethics of care will find a genuine tending of persons; political economy will find an extraction; feminism will divide against itself. That they do not agree is not a defect of the inquiry to be repaired by a final synthesis. It is, I will propose, what \textit{amai} has to teach. The goodness or badness of a cycle that looks good is not the kind of thing that lives inside a single framework's jurisdiction. It lives in the polyphony, and the honest form for such a truth is a set of conclusions in the plural, with no one of them permitted the last word.

\section{The Missing Mediation: Lack, the Cute, and Amai}

To ask what is exchanged in a sweet cycle, one must first ask what \textit{amai} is, and the answer requires a theoretical foundation that this paper, rather than gesturing toward in another text, will lay out in full here. The foundation is psychoanalytic, and it concerns the way a symbolic order generates, sustains, and mediates a lack around which subjects and their desire are organized. Because the argument is the structural ground for everything that follows, on the cute, on \textit{amai}, on the political economy of the sweet cycle, and on the practices that close the paper, it is given its own section and developed without abbreviation. I have presented a version of this argument elsewhere in Chinese; what follows is its full development in English, recast so that it flows into the analysis of \textit{amai} that is this paper's concern.

\subsection{A Note on the Status of the Concepts}

A clarification must precede the argument, because the argument uses terms that are easily and damagingly misread. The Lacanian vocabulary that follows, and in particular the expressions ``feminine position,'' ``masculine position,'' and ``hysteric structure,'' are technical terms internal to a particular psychoanalytic theory. They denote \textit{structural positions within a symbolic system}, not empirical claims about women, men, or any actual person. The ``feminine position'' in Lacanian theory names a structural relation to the symbolic order, the position of what cannot be fully inscribed within the order's central organizing logic, and it can be occupied by a subject of any gender. It is not a statement about biological sex, social gender, personality, or the psychology of real individuals. Likewise the term ``hysteric structure'' is, in this theory, a technical name for a particular mode of relation between subject and symbolic order, and bears no relation to the term's colloquial or pathologizing uses.

This matters for two reasons. First, intellectual honesty requires acknowledging that the theory reconstructed below was produced in mid-twentieth-century European psychoanalysis and rests on assumptions, the universality of a certain symbolic structure, the centrality of a single organizing signifier, the universality of the Oedipal narrative, that have been extensively criticized by feminist theory, queer theory, and post-structuralism, and that this paper does not endorse as descriptions of how gender or subjectivity must be. The aim in reconstructing the theory is to understand its internal logic and to extract from it a structural insight about lack and mediation that the later sections need, not to defend its presuppositions or to apply them as norms to actual people. Second, the later feminist section of this paper will itself turn a critical eye on exactly these presuppositions, and the reader should hold the reconstruction that follows as a theoretical instrument under examination, not as an endorsed account of women. With that frame in place, the argument can be given its full strength, which is what it requires in order to do its later work.

\subsection{The Problem of the Feminine in Psychoanalytic Theory}

Psychoanalysis confronted a theoretical difficulty about the feminine from its beginning. Freud confessed his perplexity before feminine psychology in the famous question of what a woman wants, and the question marks a structural obstacle the early theory met rather than a mere gap in observation. Freud tried to understand the formation of the feminine subject through a series of concepts concerning sexual development, but his frameworks took male development as their reference and placed the feminine in a derivative, supplementary position. This difficulty exposed not only the limits of early psychoanalysis but, more deeply, a fundamental problem in the theoretical system of the time when it came to treating the feminine position.

It was not until Lacan's work in the middle decades of the twentieth century that psychoanalysis advanced a new understanding of the feminine position. By importing structural linguistics and Saussurean semiotics, Lacan re-grounded psychoanalysis upon an analysis of the symbolic order. Within this frame he advanced the well-known assertion that woman does not exist, \textit{La femme n'existe pas}. The theoretical content of the claim is this: at the level of the symbolic system, there is no fully definable, fully nameable feminine essence.

The claim bears on how the subject is constituted by desire, how it enters the order of language, and how it relates to the desire of the other. For Lacan the symbolic order is a system built around a central signifier, for which he used the term \textit{phallus}, understood here as a signifier of symbolic function derived from his symbolic reworking of the Oedipus complex, and it organizes subject positions through the presence and absence of that signifier. In Lacan's theoretical assumption, a subject occupying the masculine position can obtain a stable location within this system of having and not-having by identifying with the symbolic position of the father, whereas a subject occupying the feminine position remains outside that logic. ``Outside'' here means not fully inscribable within the symbolic system. The feminine position is, in Lacan's account, always in excess, remaining, not fully nameable, and this is precisely what constitutes a particular structure of subjectivity, which in the psychoanalytic tradition has been called the hysteric subject, again a strictly technical designation that has nothing to do with the colloquial or pathological sense of the word.

When Lacan discusses the association of the feminine position with this particular structure of subjectivity, what concerns him is an isomorphism between a subjective structure and the symbolic order, well beyond any clinical symptom or pathology. In his theoretical system the feminine position and this subjective structure exhibit a correspondence within the symbolic, because both confront the same fundamental predicament: how to come into being and sustain subjectivity within an order that cannot provide them a complete name. The predicament, on Lacan's analysis, arises from a structural asymmetry of the symbolic order itself rather than from any essential difference between the sexes.

\subsection{The Self-Reference of the Symbolic and the Emergence of Lack}

To understand how this mechanism is generated, one must return to its ontological presupposition. In Lacan's framework the operation of the symbolic order depends on a self-referential structure. The generativity of the cosmos, the continuation of life, the reproduction of society, and desire itself all must point back toward themselves in order to reproduce. This self-reference points toward the very possibility of continuation, beyond any external terminus. The symbolic order exists, on this account, so that being may go on being generable, not in order to finish saying being.

Yet a core thesis of the theory is that the symbol can never exhaust the real. This is the very condition under which the symbolic order keeps operating. The symbol must designate being, and yet it necessarily produces, within itself, a point that cannot be said in full. If lack were to appear directly, Lacan holds, desire would collapse, meaning would break off, and the subject would slide toward a psychotic rupture. The symbolic order therefore faces a structural requirement: lack must be seen but not faced, sensed but not possessed.

From this arises the necessity of one of the theory's central concepts, the \textit{objet petit a}. It is a mediating position that neither coincides with lack nor fills it. The \textit{objet a} functions as the cause of desire, \textit{cause du désir}, rather than its satisfaction. In Lacan's construction, lack emerges structurally as the condition that makes self-reference possible. Any symbolic society, when it runs stably, will almost necessarily generate this structural solution, because a society must continuously reproduce desire, or else the subject will cease to act, to connect, to create. Desire, in Lacan's theory, must be sustained through mediation, in a process that can be completed neither by satisfaction nor by cancellation.

\subsection{Paths of Subject Formation}

In Lacan's framework the formation of a subject occupying the masculine position is described thus: through the Oedipus complex and the intervention of the name of the father, the subject undergoes a process of symbolic identification. The boy, by identifying with the father, obtains within the symbolic order a position associated with the symbolic position the father is supposed to hold. Lacan abstracts this mechanism into a distinction of having and not-having organized around a central signifier: the father is supposed to have this signifier, the mother is positioned as lacking it, and the masculine subject, by identifying with the father, obtains a symbolic location on the side of having. In this schema subjectivity amounts to obtaining a determinable position within the system of the distinction, a closed process that can be completed by symbolic relations alone.

It must be said that this binary distinction is Lacan's theoretical construction after a symbolic reworking of the Freudian Oedipus complex, and that the construction rests on an assumed universality of the Oedipus complex. The symbolic system need not take the form of this binary. The assumption is itself widely contested in contemporary theory, a point the later feminist section will press.

Lacan's theoretical observation is that the symbolic order can hardly furnish the subject of the feminine position with a closable place. In his assumption, because the symbolic order is organized around the central signifier and its distinction of having and not-having, and because in the classic Oedipal narrative the feminine is positioned as the one lacking the signifier, the feminine position cannot be fully mapped into the system of distinction. In Lacan's description the feminine position lies in an unclosed zone between symbol and real. This position can hardly complete a subjective closure by symbolic relations alone, and so Lacan observes that the subject of this position tends to seek, test, and generate its own position within the structures of relation in the register of the real.

Lacan uses the concept of the not-all, \textit{pas-toute}, to describe the relation between the feminine position and the symbolic order. The not-all names a fundamental condition of not being wholly subsumable under symbolic logic, and exceeds the senses of partial or incomplete. This construction attempts to explain certain of Lacan's clinical observations, observations that, it must again be stressed, were grounded in the clinical practice and social context of the mid-twentieth century and whose universality and applicability are widely questioned today.

In Lacan's schema these phenomena are read as one possible outcome of a path of subject formation. When the symbolic system can hardly supply a completed position, the theory holds, the subject may seek the confirmation of its existence more in the continual generation of relation. But the subject is a relational subject, relation requires a continuous dynamics to sustain it, and the core of that dynamics in Lacan's theory is precisely desire.

\subsection{From Structure to Generation: The Missing Mediation}

Lacan's theory offers not only a static structural analysis but an account of the subject's becoming. Turning to the generative perspective brings the theory's key insight into view. The crux is the dynamics of relational structure and its sustainability. In Lacan's frame, a possible difficulty the subject of the feminine position faces is this: standing in a position to which the order can hardly furnish an adequate name, how is relation to be kept running?

Lacan's analysis holds that any relation, to continue, may need to rely on desire, on tension, and on incompleteness. His theory of desire states that the structure of desire points toward the continuation of lack, beyond pointing toward any object as such. One may render the theoretical premise as a chain of dependencies: in Lacan's frame, the stability of the subject is bound up with the continuation of relational dynamics, relational dynamics with the continuation of desire, and desire with the continuation of lack, so that the subject's stability is, in a degree, bound up with the persistence of lack. In Lacan's theory, lack can become a bearer of stability.

But Lacan analyzes further that if lack is merely an external condition, then once lack is filled, displaced, or denied by the other, the dynamics of relation may break off. For a subject position more dependent on relational structure, this may constitute a certain risk. Hence a generative mechanism the theory proposes: to let ``my own self'' become the site at which lack can appear, be approached, and be sustained, so that lack is transformed from an external condition into an internal subjective function. This process may convert the persistence of lack into a function of the subject itself, which is one theoretical description of the position of the \textit{objet a}.

When the subject ties its stability more to the persistence of lack, Lacan's theory conjectures, it may, at the level of experience, present in certain patterns, confirming through being desired, being needed, being aimed at, that relation is still being generated. In Lacan's theory, to be desired by the other can become a sign that lack is still in operation.

It is at exactly this point that an extension of the theory toward cultural forms becomes possible, and it is the hinge on which this paper turns from psychoanalysis toward \textit{amai}. Some contemporary scholars, extending Lacanian theory, have analyzed the operation of particular cultural symbols. In discussing certain aesthetic categories, Sianne Ngai has indicated that the cute, \textit{kawaii}, as a cultural representation, may operate differently from the beautiful. If the beautiful inclines toward completion, symmetry, and stability, the cute may instead contain incompletion, a slight asymmetry, and a summons toward response and attention. From the standpoint of Lacanian theory, such a cultural symbol may operate through the not-all, through marginality, through unclosedness, and so play a particular role in the economy of desire.

It must be stressed that this kind of cultural analysis depends heavily on a specific theoretical frame and cultural context. Any attempt to apply a theoretical model to concrete cultural practice demands great caution, because the meaning of a cultural symbol is plural, fluid, and situation-dependent. The discussion here serves as one possible extension of Lacanian theory, and should not be read as a normative interpretation of any cultural phenomenon, still less as an evaluation of any group's behavior or aesthetic choices. In the theoretical language of Lacan, certain cultural representations may function as a mediation of lack, but the relation between this theoretical description and concrete cultural practice and individual experience always requires a critical distance and an openness to reflection.

\subsection{The Internal Logic of the Schema}

To gather the reconstruction: because the subject of the feminine position can hardly complete a subjective closure through the symbolic order, its stability may depend more on the continual generation of relation in the register of the real. Under the assumption that the continual generation of relation depends on desire, and that desire operates on the condition that lack persists, Lacan's analysis proposes that, in order to sustain its stability, the subject of the feminine position may tend to embed the persistence of lack as a part of its own subjective function, so that lack is transformed from a fillable external condition into an internal component of the subject. Under this generative perspective, Lacan's analysis holds that the feminine position may incline toward the position of a mediation of lack, and may present, at the level of experience, as a particular structural pattern of subjectivity.

The core of this structure, in Lacan's account, is to take oneself as the site at which the desire of the other may be continually kept in question. The subject of the feminine position, because the symbol can hardly close, may need to generate itself within relation, and because relation may need to be sustained through desire, may incline to take the question of what I am in your desire, the \textit{Che vuoi?}, as a central problem through which the subject continually generates itself. This mechanism, Lacan suggests, may arise under the structural condition that, in a society centered on the symbol and operating through the maintenance of desire, lack needs to be made accessible without being eliminated. Historically and structurally, the feminine position may tend to occupy the place of the mediation of lack, which may be a result of the symbolic system's search for stability.

\subsection{Building on the Reconstruction: How the Cute and Amai Become Legible}

The reconstruction has done its structural work, and the paper can now draw from it the consequences that the analysis of \textit{amai} requires. The first consequence concerns the cute. If the cute operates, as the extension above proposes, through incompletion and the summons to response, through holding lack open rather than closing it, then the cute is not merely a taste or a style but a form with a precise function in the economy of desire: it is a culturally elaborated mediation of lack. The cute object does not satisfy; it solicits. It keeps the question of the other's response alive by presenting itself as unfinished, as needing to be answered. This is why the cute can sustain an endless relation to itself where the beautiful, inclining toward completion, tends to release the beholder into repose. The cute does not release; it calls again.

The second consequence brings us to \textit{amai} itself, and here the reconstruction pays its full debt. \textit{Amai}, the licensed reliance through which a bond knows itself, can now be seen as a relational enactment of exactly the structure the reconstruction describes. \textit{Amai}, the adjectival root of \cn{甘え}, names in the Japanese psychological tradition, in Takeo Doi's classic treatment, the presumption upon another's indulgence. What the one who is indulged seeks is not, at bottom, the object through which the indulgence passes; it is the confirmation of being wanted, the sign that lack is still in operation between the two, that the other's desire still turns toward oneself. In the vocabulary the reconstruction supplies, the one who practices \textit{amai} takes up, within the relation, the position of the mediation of lack, making of the bond a site at which the question \textit{what am I in your desire} can be continually posed and continually, partially, answered. Doi's account and Lacan's reconstruction, arrived at from entirely different traditions, here describe one and the same relational form.

Doi's analysis deserves to be held a moment longer, because it adds a feature the reconstruction does not by itself supply and the later sections will need. For Doi, \cn{甘え} is not a private feeling but a relational presupposition, the often unspoken expectation that the other will receive one's reliance with goodwill, and it is constitutive rather than incidental: a relation in which such reliance could not be presumed would not, on this account, be an intimate relation at all. Two consequences follow. First, \cn{甘え} is structurally reciprocal even when enacted asymmetrically, for the one who presumes and the one who indulges are not opposed roles but two moments of a single relational form that requires both. Second, and this is the feature the critical sections will exploit, the presumption can be true or false to itself. It can be met with the goodwill it presumes, in which case the reliance is answered and the bond is real; or it can be met with a simulation of goodwill, a performance of indulgence that is in fact extracting something, in which case the very structure that makes intimacy possible becomes the structure through which intimacy is counterfeited. \textit{Amai} is thus, from the start, a form that can be inhabited well or ill, and Doi's account already contains, without naming it, the possibility of a forged version of the thing.

The third consequence is the one the whole later argument turns on. If \textit{amai} is the relational mediation of lack, then the objects of the sweet cycle, the cosmetics, the small accessories, the things gathered under the sign of sweetness, are the material mediations of that mediation. They are the \textit{objet a} in commodity form. To purchase the sweet thing is to purchase access to the question, \textit{am I still wanted}, and to receive, in the partner's indulgence, the sign that the question is still alive. And because the cute object holds lack open rather than closing it, the purchasing does not terminate. A satisfied desire would end it; a desire mediated by the cute is, by design, never quite satisfied, and so it returns. This is the precise structural reason the sweet cycle can renew itself without end, and it is the reason, the political economy will argue, that capital found in \textit{amai} so apt a structure to occupy.

\subsection{Need, Demand, Desire: The Triad the Later Analysis Requires}

One distinction internal to the reconstruction must be set out plainly on its own, because the whole later analysis of how capital operates on \textit{amai} depends on it. A need, in this vocabulary, is a deficit that an object can fill: hunger is filled by food and then is gone until it returns. A demand is what a need becomes when it must pass through another person and through language to be met, and because it passes through another, a demand is never only a demand for the thing, it is always also a demand to be the kind of being whose demand the other will answer, a demand, that is, for love. Desire is the remainder that no satisfaction of demand exhausts, the surplus of the demand for love over any particular thing given in response. This remainder is structural and permanent. One can be given the thing one asked for and still not have been given the recognition one was asking through it, and the gap between the two is where desire lives. The reason this matters for \textit{amai} is that the cycle operates precisely on the gap. It offers, through the commodity, an endless series of partial answers to a demand that by its nature cannot be finally answered, and the endlessness is not a malfunction of the cycle but its engine. A cycle built on need would terminate when the need was met. A cycle built on desire renews itself forever, because the remainder it addresses is never used up.

It is essential, and the later critical sections will turn on this, that none of this is in itself sinister. That a relation runs on a lack kept open is not a flaw in the relation; it is, on the reconstruction's account, what makes it a living relation rather than a closed and finished arrangement between two satisfied parties who no longer have any question to put to one another. The danger enters not with the structure of desire but with what a particular social organization does to that structure, and the central claim of the political economy to come is that capital has learned to occupy the gap, to insert itself as the necessary supplier of the partial answers, and to enlarge the gap deliberately so that ever more answers must be bought. At this stage the point is only that \textit{amai}, read through need, demand, and desire, is a structure whose very openness, the openness that makes it intimate and alive, is also the opening through which it can be colonized.

\subsection{The Image of the Sweet Self}

There is a further layer, and it is the layer at which a later section's question about flourishing will be decided, so it is named here and left armed. What is consumed in the sweet cycle is not only an object and not only the partner's indulgence. It is an image of the self as the one who is indulged, the sweet self, the self worth being charmed by. This is consumption in the register Baudrillard described, where what circulates is the sign and not the use, and it is capture in the register the present series calls the imaginary, where a living dynamic is fixed into a flattering picture of itself. The danger is not that the image is false. The danger is that the image can substitute for the very thing it pictures. A self can come to consume the image of being indulged in place of the slower, riskier, genuinely generative work of unfolding its own singular dynamic with another. Whether the sweet cycle opens that unfolding or forecloses it, whether the image is a medium of generation or a cage that stands in for it, is precisely the question this paper will later put under the name of the eudaimonic.

It is worth being precise about why the imaginary capture is so difficult to detect from within, since the difficulty is what gives the forged cycle its cover. The image of the sweet self is not experienced as an image. It is experienced as oneself, as who one simply is, and the labor of producing and maintaining it disappears into the feeling of merely being oneself well. This is the imaginary's characteristic operation: it does not announce itself as a representation, it presents itself as immediate reality. A person caught in the image does not feel they are consuming a picture; they feel they are being authentic, that the small bright thing expresses rather than constructs them. And because the partner too responds to the image as though to the person, the cycle can run for a long time on the smooth surface of the image without anyone touching the singular dynamic underneath. Nothing in the experience flags the substitution. This is why the eudaimonic criterion, when it arrives, cannot be applied by asking the participants whether they feel fulfilled; they may feel entirely fulfilled by the image precisely in the measure that the image has succeeded. The criterion will have to ask a harder question, about whether the relation is generating anything genuinely new between two unfolding singularities, or only polishing two pictures to a higher shine.

\subsection{What the Foundation Opens: Toward Political Economy and Practice}

Two contemporary considerations extend the reconstruction in the direction the rest of the paper will travel, and naming them here closes this section by handing its insight forward. The first is that the stability of the symbolic order which Lacan presupposed is itself, in the present, under new strain. In an age of highly symbolic technologies, of large language models, algorithmic decision systems, and automated symbolic production, the generation, distribution, and reproduction of meaning are increasingly detached from a human, subject-centered structure. This re-poses the question of how the subject is constituted, and it makes the relational, affective, and ethical dimensions of existence newly salient. It is against this background that feminist theory, the ethics of care, and the discussion of relational capacities such as empathy and co-presence offer crucial resources, resources that think the subject from the standpoint of relational being and that exceed any single question of gender. The reconstruction's value, in this light, is not as a fixed account of women but as a historicized lens on how the relational capacities long positioned at the margin, the capacities of the mediation of lack, come to the center precisely when the symbolic order trembles.

The second consideration is the one that drives the remainder of this paper. If the mediation of lack is a structural position that a symbolic society generates because it must continually reproduce desire, then a society organized by capital will not leave that position unoccupied. It will discover, occupy, and industrialize it. The reconstruction shows that \textit{amai} is the relational mediation of lack and that the cute object is its material form. The political economy to come will show that capital has learned to insert its commodity into exactly the place where the mediation of lack does its work, to supply the partial answers the demand for love requires, and to enlarge the lack deliberately so that the supply must be bought again and again. The structure the reconstruction describes as the condition of a living relation becomes, under capital, the structure of an extraction, and the practices that close the paper will ask whether the living relation can be recovered from the extraction that has colonized it. The foundation laid in this section is thus not a digression into psychoanalysis but the precise diagram of the field on which the rest of the argument is fought. The structure of \textit{amai} is that field. The question the following sections take up is who is farming it, at whose expense, and how it might be farmed otherwise.

\section{The Strongest Version of the Good Cycle}

Before the cycle is tested it must be granted its full strength, and granting it is not a rhetorical concession but a methodological requirement. A critique that defeats a weak version of its object has defeated nothing. The sweet cycle deserves its strongest formulation, and that formulation is available in the language of generative justice, where an arrangement counts as good when it generates value and returns that value to those who generate it, rather than siphoning it to a party who contributes nothing to the generation. Read within its own boundary, the sweet cycle satisfies this standard with a completeness that should give any critic pause.

The standard deserves to be stated carefully, because it is the most demanding form of the case for the cycle and the form the rest of the paper must contend with. Generative justice, as developed in the work this series engages most closely, departs from the familiar liberal framing in which justice is a matter of distributing a fixed stock of goods fairly after they have been produced. Its concern is upstream of distribution. It asks how value is generated in the first place and whether the value returns to those whose activity generated it, or is alienated from them and accumulated elsewhere. The paradigm of injustice on this account is not an unequal slice of the pie but an arrangement in which those who bake are systematically separated from what they bake, so that value flows away from its generators to a point of concentration. The paradigm of justice, correspondingly, is a cycle in which the value generated by an activity returns to nourish the activity and those who carry it out, a self-sustaining loop with no leak to an external accumulator. This is a strong and attractive ideal, and the reason the sweet cycle is such a formidable case is that, viewed within its own boundary, it appears to instantiate the ideal almost perfectly.

Trace it once more, slowly, granting each step its due. A firm offers an object in the register of the cute. A person acquires it and is thereby enabled to bring a certain sweetness into a relationship. The partner receives this sweetness as an address and answers it with indulgence, and in answering feels the specific reward of being the one who is relied upon, the one whose care is wanted. The first person, indulged, receives the confirmation they sought, that they are wanted, that the lack between the two is still in living operation. The bond is, by this exchange, not merely maintained but deepened, because each party has given the other something the other genuinely values and has received something in return. The deepened bond raises, rather than lowers, the willingness of both to invest again in the small mediations that let the exchange recur. The firm, profiting, is enabled to continue offering. Around the loop, value has been generated, the value of a strengthened intimate bond, and it has been returned to those who generated it. No party named in the loop is a mere instrument of another. Each is at once a contributor and a beneficiary.

One should resist the temptation to weaken this picture prematurely by pointing out that the firm's motive is profit. The objection is too quick, and answering it strengthens the case rather than undermining it. That a party acts from interest does not, on the generative account, make an arrangement unjust; the baker who is nourished by the bread acts from interest too. What would make it unjust is a leak, a flow of value out of the loop to someone who contributes nothing to the generation and returns nothing to it. Within the boundary as the strongest version draws it, there is no such leak. The firm contributes the object that mediates the recognition and is paid for that contribution; its profit is, on this reading, its share of a value it helped generate, not a siphoning of value it did not. So long as the boundary holds where the strongest version places it, the profit motive is not an embarrassment to the cycle's justice but is simply absorbed into it. The cycle does not merely tolerate the firm; it justifies the firm's return. This is how complete the strongest version is, and it is why the critique to come cannot proceed by moral disapproval of commerce. It must proceed by challenging the boundary itself.

One can put the point in the vocabulary the series has developed and the case only grows stronger. The cycle does not look like sterile repetition, the mere going-around that returns each party to exactly where they began with nothing accumulated. It looks like a spiral. Each turn deepens the bond a little further than the last; something is accumulated that was not there before; the return is not to the same point but to a higher one. In the series' terms the loop appears to carry a positive holonomy, a real sublimation of value across the cycle rather than a flat circulation. And it appears to carry that holonomy internally, generated by the participants' own exchange, not imported from some reservoir outside. If that appearance held, the cycle would be not merely permissible but exemplary, a small working model of how intimate life and economic life might be made to reinforce one another without either corrupting the other.

I want to leave this section without a single qualification attached to it, because the qualifications are the work of the sections that follow and they should not be allowed to leak backward and weaken what must first stand at full height. The strongest version of the good cycle is genuinely strong. Anyone who feels its force is feeling something real. The discomfort that the Ginza cover also provokes does not come from any weakness in this account. It comes from a suspicion that this account, however true within its frame, has drawn its frame too small, and that the goodness it correctly perceives is purchased by not looking past the edge. To look past the edge is the task that begins now. But the account just given is the thing that must be looked past, and it is not a thing to be ashamed of having found compelling.

\section{Many Frameworks, One Cycle: Incommensurable Readings}

The strongest version having been granted, the cycle can now be approached from several directions at once. This section places six frameworks before it, psychoanalysis, neuroscience, the ethics of care, structuralism and linguistics, dialectics, and Kantian ethics, and asks each not only what it sees but what it cannot see. The two heaviest frameworks, political economy and feminism, are held for sections of their own, and the formal instruments that can be brought to bear, the dynamical, grammatical, field-theoretic, categorial, game-theoretic, and information-theoretic, are gathered in the section that follows this one, since they are tools of a different kind. The discipline observed throughout is that no framework is allowed to deliver the others' verdict. Each illuminates one face of \textit{amai} and goes dark at the edge of its competence, and the darkness is reported as carefully as the light.

A word is owed on why the frameworks are kept incommensurable rather than integrated, since the reflex of much scholarship is to seek a synthesis in which each contributes a piece to a single completed picture. The refusal of synthesis here is not eclecticism or a failure of nerve. It rests on a claim about the object: that the goodness of a cycle which looks good is not a unitary property awaiting measurement by a sufficiently comprehensive instrument, but a question that genuinely receives different and irreconcilable answers depending on the scale and the stance from which it is asked. Psychoanalysis asks at the scale of the desiring subject and answers in the currency of desire. Neuroscience asks at the scale of the nervous system and answers in the currency of neural event. The ethics of care asks at the scale of the face-to-face relation and answers in the currency of responsiveness. Structuralism asks at the scale of the sign system and answers in the currency of difference. Dialectics asks at the scale of the self-moving process and answers in the currency of contradiction. Kantian ethics asks at the scale of the rational will and answers in the currency of the maxim. These are not several views of one thing that could be added together; they are several things, each real, that the single word \textit{amai} happens to name at once. To force them into a synthesis would be to choose, covertly, one framework's currency as the master currency into which the others must be converted, and that covert choice is exactly the move the paper's method is designed to refuse. The honest alternative is to let each speak fully and to let the dissonance stand as a finding.

\subsection{The Epistemology of Psychoanalysis}

Psychoanalysis sees the sweet cycle better than any other framework at one specific thing: it sees why the cycle never ends. The earlier account of \textit{amai} as a demand for love mediated by the \textit{objet a} explains what no account of pleasure or utility can, namely that the satisfaction the cycle delivers is structurally incomplete and must be incomplete in order to recur. If the purchasing satisfied a need, the purchasing would stop. It continues because what is sought is not a need's satisfaction but the maintenance of a desire, the standing confirmation that one is still wanted, and desire is sustained only by a lack kept in operation. Psychoanalysis can therefore say something true and non-obvious: the endlessness of the cycle is not a sign of addiction or of manipulation alone, it is the signature of desire as such, and a relation that kept desire alive would also, for that very reason, never be done.

But this is also the exact place where psychoanalysis must report its own limit, and the report is part of its contribution rather than an embarrassment to it. Psychoanalysis can tell us that desire is in operation. It cannot tell us whether the cycle in which that desire operates is just. The structure of demand and lack is the same whether the loop is internally generative or sustained by extraction from outside; the \textit{objet a} does its work identically in a cycle that wrongs no one and in a cycle that runs on hidden labor. The framework that best explains why the cycle persists is constitutively silent on whether the cycle is good. This silence is not a gap to be filled by more psychoanalysis. It is a boundary of jurisdiction, and noting where it falls tells us that the question of justice will have to be carried by frameworks of a different kind.

There is a further and subtler contribution that psychoanalysis makes precisely by way of its silence, and it concerns the seductiveness of the cycle to thought itself. Because the structure of desire is so satisfying to describe, so elegant in its account of why the cycle renews, there is a temptation to feel that in having explained the cycle one has also exonerated it, that to show how natural and structural the endlessness is, is to show that there is nothing to object to. This is a false inference and a dangerous one, and psychoanalysis, properly understood, warns against it. To explain a mechanism is not to license it. The same desire that makes intimacy alive can be harnessed to an extraction, and the elegance of the explanation of desire says nothing about whether the particular harness is just. Psychoanalysis thus offers not only an account of the cycle but a caution against its own seductiveness as an account: it can make the cycle intelligible without making it innocent, and the reader who feels the explanation settle their unease should treat that settling as itself a datum to be examined rather than a verdict to be trusted. The framework that explains desire is, in this respect, the framework most likely to be mistaken for an absolution, and it disclaims the role explicitly.

\subsection{Neuroscience: The Reward and Attachment Substrate of Amai}

Neuroscience contributes something the more interpretive frameworks cannot, an anchor in the body that prevents the whole discussion from floating off into pure construction. The satisfactions of the sweet cycle are not only meanings; they are events in the nervous system. The indulgence at the heart of \textit{amai} engages the circuitry of attachment and reward. The bonding that follows shared care recruits the neuropeptide systems, oxytocin among them, that subserve pair bonding and social trust, and the small recurring gratifications of the cycle engage the dopaminergic reward pathways that track not consummation so much as anticipation, the wanting that precedes and outlasts any particular getting. The objects of the cycle are largely objects in the register of the cute, and the cute is not an arbitrary cultural taste. The infantile schema, the large eyes and rounded forms that Konrad Lorenz first described and that later work has tied to activation in reward-related regions on a rapid timescale, recruits a caregiving and reward response that is to a significant degree pre-cultural. When the Ginza cover trades in cuteness it is trading, in part, on a substrate the nervous system brought to the encounter.

What neuroscience establishes is therefore the reality of the satisfaction, and this strengthens, rather than weakens, the strongest version of the good cycle. The participants are not deluded; their reward systems are genuinely engaged; the bond that forms is underwritten by the same machinery that underwrites attachments no one would dream of calling unjust. But the limit of the framework is as sharp as its contribution, and it is the most important negative result in this paper. Neural reality is not justice. A cycle can be completely real in the reward and attachment systems of everyone inside it and still be forged at the level of structure, sustained by an extraction those reward systems have no way of registering. The nervous system reports that the satisfaction is occurring. It does not and cannot report where the conditions of that satisfaction were manufactured or at whose expense. Any argument that moves from \textit{they are genuinely happy} to \textit{the cycle is therefore good} commits precisely the error this subsection forbids. The happiness is genuine. The inference is not licensed.

It is worth registering, without overstating, the temporal structure that the reward literature suggests, because it bears directly on why the cycle renews rather than terminates. The dopaminergic system, on the dominant reading, tracks not the receipt of a reward so much as the prediction and the anticipation of it, the gap between expectation and outcome, with the strongest signal at the moment of anticipated rather than consummated reward. If this is even approximately right, then the neural substrate is not a substrate of satisfaction in the sense of a deficit filled and quieted, but a substrate of wanting, of the lean toward a reward not yet had. This is a striking convergence with the psychoanalytic account, arrived at from an entirely independent direction and in an entirely different vocabulary: the structure that keeps the cycle turning is, at the level of the reward system as at the level of desire, a structure of the not-yet, of an anticipation that the having does not extinguish. I draw the convergence cautiously, as a resonance between two frameworks and not as a reduction of either to the other, which would violate the method. But the resonance is instructive. It suggests that the endlessness of the sweet cycle is overdetermined, written into the subject at more than one level, and therefore that any practice hoping to make the cycle good cannot proceed by trying to abolish the wanting, which is constitutive, but only by altering what the wanting is harnessed to and at whose expense it is fed.

The limit must then be stated once more in its sharpest form, because it is the hinge on which the paper's later refusal of an easy verdict turns. That a satisfaction is real in the nervous system, that it is overdetermined, ancient, and shared with attachments beyond reproach, settles nothing whatever about justice. A forged cycle does not feel forged. It feels, to the reward and attachment systems of those inside it, exactly as a genuine cycle feels, because those systems are reporting on the satisfaction and not on its conditions, and the conditions are precisely what the forgery hides. This is why no amount of evidence that participants flourish subjectively can establish that a cycle is just, and why the appeal to felt happiness, however sincerely advanced, is structurally incapable of answering the question the Ginza cover raised. Neuroscience, having confirmed the happiness with more authority than any other framework, is for that very reason the framework that most firmly closes off happiness as a route to the verdict.

\subsection{The Ethics of Care}

The ethics of care arrives at the sweet cycle from a direction that both psychoanalysis and political economy tend to foreclose, and it begins by refusing a prejudice the others smuggle in. In the tradition that runs through Carol Gilligan and Nel Noddings, dependence is not the opposite of maturity and not a deficiency to be outgrown. The autonomous, self-sufficient agent of liberal moral theory is, on this view, a fiction that has systematically devalued the relational capacities, receptivity, responsiveness, the willingness to be relied upon and to rely, on which every actual human life depends. Read through this lens, \textit{amai} is not a regression to be diagnosed but a competence to be recognized. To be able to presume upon another's indulgence, and to be able to welcome that presumption, is to be skilled in a grammar of relation that a culture of self-sufficiency has taught us to disparage. The care ethicist can therefore see what the political economist, intent on the hidden labor, is liable to miss: that inside the cycle, persons may be genuinely tending one another, answering one another's address, treating one another as ends.

The framework's question is accordingly its own, neither the analyst's nor the economist's. It asks whether, within the relation, persons are met and responded to, or merely used. And it can return a verdict that the other frameworks cannot reach: that a sweet cycle may be, in the conduct of the parties to one another, ethically nourishing, a real practice of mutual care, even as the same cycle is something else again when viewed from the standpoint of what sustains it. This is the first appearance of a fracture that will organize much of what follows, the possibility that a cycle is ethically good in the treatment of those within it and unjust in the structure that feeds it. The ethics of care holds one side of that fracture open and insists it not be closed prematurely. What it cannot do, and does not pretend to do, is see across the relation to the chain of production and reproduction that supplies the relation its means. For that we need the framework that looks hardest at exactly what care, attending to the face before it, tends not to see.

The fracture deserves to be named precisely, because it is easy to mistake for a contradiction to be resolved when it is in fact a feature to be preserved. It is entirely coherent, and indeed common, for a relation to be ethically good in the care of its members and unjust in its structural conditions, and the two assessments do not cancel because they are assessments of different things. The ethics of care assesses the relation as a relation: are these two people responsive to one another, do they treat one another as ends, is there genuine tending here. Justice, as the later sections will deploy it, assesses the relation as a node in a wider system: what sustains this tending, who bears its costs, where does its surplus come from. A relation can score high on the first and low on the second, and when it does, the right description is not that the care was illusory but that the care was real and was, at the same time, underwritten by an injustice the carers did not see and did not choose. This is not a comfortable result, because it denies us the satisfaction of a single coherent judgment, but it is the truthful one, and the ethics of care earns its place in the inquiry precisely by defending the reality of the care against a critique that, left unchecked, would explain it away. The economist is right that there is an extraction. The care ethicist is right that there is genuine care. The mistake would be to think that one of them must therefore be wrong.

\subsection{Structuralism and Linguistics}

The structuralist tradition, descending from Saussure through Lévi-Strauss and Barthes, approaches \textit{amai} not as a feeling or a transaction but as a position within a system of signs, and it sees, more clearly than any other framework, that sweetness has no meaning in itself but only by its difference from what it is not. To be cute is to occupy a marked term in an opposition: cute against cool, soft against hard, the one who is indulged against the one who provides. The structuralist insight is that the sweet object signifies nothing intrinsically; it signifies sweetness only within a code that assigns it that value by contrast, and the code, not the object, is the bearer of meaning. This is why the same gesture can read as charming in one system and as cloying or manipulative in another, and why the magazine cover can confer a meaning, justice, upon sweetness simply by placing the two terms in apposition within a signifying frame. Barthes would recognize the cover at once as a myth in his technical sense, a second-order signification in which a sign already complete, sweetness, is made the signifier of a further concept, justice, so that a cultural and contestable association is naturalized into an apparent identity. The structuralist can thus say something the others cannot: that the cover's proposition operates by enlisting a whole differential system, and that its persuasiveness is the persuasiveness of a code that has made its contingent associations feel like nature.

But the structuralist framework reaches its limit at exactly the point the inquiry most needs, and it reaches it for a principled reason. A system of differences has no normative dimension internal to it. Structuralism can map how sweetness comes to mean, can lay bare the oppositions and the mythological operation, but it cannot say whether the system is just, because justice is not a relation of difference within a code but a relation among persons that the code may serve or betray. The structuralist brackets the referent and the subject by method, attending only to the play of signs, and in bracketing them it brackets precisely the labor, the bodies, and the costs that the question of justice must count. Structuralism shows the cover to be a myth; it cannot, from its own resources, say that the myth conceals an extraction, because the concealed is by definition what the system of signs does not contain. It hands the inquiry a precise account of how meaning is made and a principled silence about whether the meaning is paid for honestly.

\subsection{Dialectics}

A dialectical reading, in the line of Hegel and Marx, sees in the sweet cycle neither a stable good nor a simple fraud but a process driven by its own internal contradiction. The cycle is not one thing that is then judged; it is a movement in which the good and the forged are two moments of a single unfolding. The very feature that makes the cycle nourishing, that it keeps lack in operation and so keeps the relation alive, is the same feature that opens it to colonization, and the dialectician insists that these are not two facts but one, the positive and negative faces of a single determination. The recognition that deepens the bond and the consumption that deepens the dependence advance together; the cycle develops by intensifying the contradiction it carries, until what presented itself as pure mutual flourishing reveals the extraction it always contained, and the revelation is not the arrival of new information but the working-out of what was implicit from the start. Dialectics can therefore say what the static frameworks cannot: that the sweet cycle is not accidentally both good and forged but necessarily so, that its goodness and its forgery generate each other, and that any practice hoping to redeem it must work through the contradiction rather than wish it away.

Yet dialectics carries a danger that this paper must name rather than ignore, and the danger is peculiarly sharp here because it strikes at the paper's own method. The Hegelian form of dialectic tends toward sublation, toward an \textit{Aufhebung} in which the contradiction is resolved at a higher level and gathered into a totality that comprehends it. That movement is precisely the synthesis the paper has forsworn. A dialectics that promised to lift the sweet cycle's contradictions into a reconciling whole would be offering exactly the master framework, the single currency of comprehension, that the method refuses. So dialectics enters the inquiry under a double sign: it is indispensable for seeing that the good and the forged are internally related and self-moving, and it is to be resisted at the moment it offers to complete that insight in a totalizing synthesis. The paper takes the dialectical perception of contradiction and declines the dialectical promise of its resolution, keeping the negative, the unresolved, open. This is itself a quarrel internal to the dialectical tradition, between a Hegelian closure and the negative dialectics that, with Adorno, refuses the consoling whole, and the paper stands, here, with the refusal.

\subsection{Kantian Ethics}

Kantian ethics brings to the sweet cycle the most exacting and least sentimental of the moral questions, the question of whether anyone within it or beyond it is treated merely as a means. The formula of humanity, that one must act so as to treat humanity, whether in one's own person or in that of another, never merely as a means but always also as an end, gives a sharp instrument precisely where the ethics of care is most indulgent. Where the care ethicist asks whether persons are tended, the Kantian asks a harder question that tending alone does not answer: whether the one who is tended, and the one who tends, and above all the ones across the cycle's edge who make its objects and bear its costs, are each treated as ends in themselves or used as instruments for an end that is not theirs. By this test the inner relation may pass, for the indulged and the indulgent may well regard each other as ends. But the production worker whose labor is externalized, and the self whose unwaged self-fashioning is consumed as if it were leisure, are treated, in the cycle as organized, very largely as means, as the instruments through which others' sweetness is supplied, counted in none of the ends the cycle serves. The Kantian thus arrives, by a wholly deontological route that owes nothing to the political economy, at the same edge the economist found, and the convergence is significant: two frameworks that share no premises locate the wrong in the same place, the reduction of persons to fuel.

The limit of the Kantian framework is the obverse of its strength, and it is the same formalism that makes it powerful. The categorical imperative abstracts from all particularity to reach a universal law, and in doing so it is constitutively ill-suited to a phenomenon as culturally specific and as relationally textured as \textit{amai}. Its universality cannot easily accommodate the claim, which the ethics of care and the later cultural section both press, that dependence and the licensed reliance of \textit{amai} are goods whose meaning is internal to a particular form of relation and a particular culture, not derivable from a law that would hold for all rational beings as such. The Kantian is liable to read \textit{amai} either as morally indifferent, a mere inclination, or as suspect, a heteronomy of the will before another's desire, and in both readings it misses what the care ethicist insists is really there, a competence and a good. So Kantian ethics and the ethics of care, both moral frameworks, are themselves incommensurable, the one demanding universality and autonomy, the other defending particularity and dependence, and the sweet cycle is a case where they genuinely disagree. The paper does not adjudicate between them. It records that the deontological eye sees the instrumentalization at the edge with unmatched clarity and is, for the same reason, partly blind to the relational good at the center.

\section{The Epistemology of Formal Tools: Formalizing the Cycle and the Limits of Formalization}

The frameworks of the previous section were substantive: each delivered a reading of \textit{amai} in a currency of its own, desire, neural event, care, difference, contradiction, the maxim. This section turns to instruments of a different kind. Dynamical systems, generative grammar, field theory, category theory, game theory, and information theory do not, in the first instance, tell us whether the sweet cycle is good. They offer formal languages in which the structure of the cycle can be written down with precision. Their inclusion is not ornamental. A recurring claim of this paper is that the difference between a genuine and a forged cycle is a difference of structure, and structure is exactly what formal tools are built to capture. But the section has a second and equally important purpose, which is to show that each formal tool, in capturing the cycle, must make a choice about what to include in its representation and what to leave out, and that this choice is itself a drawing of a boundary. The formal tools are therefore not the view from nowhere that would settle from above what the substantive frameworks could not. They are themselves situated, each seeing the cycle through the particular abstraction it imposes, and the limit of each is as instructive as its reach. The reflexive lesson, stated once here and returned to at the section's close, is that to formalize is already to choose a boundary, and the choice of boundary is, as the political economy will insist, the very thing on which the judgment of the cycle turns.

\subsection{Dynamical Systems}

The most natural formalization treats the sweet cycle as a trajectory in a state space. Let the state of the relation at a given time be a point whose coordinates record the relevant quantities, the strength of the bond, the level of accumulated recognition, the intensity of felt lack, the rate of consumption, and let the cycle be the evolution of this point under a dynamics. In these terms the distinctions the paper has drawn acquire sharp images. A merely circulating cycle is a closed orbit, a limit cycle that returns to the same loop period after period with nothing gained. A genuinely generative cycle is a spiral that climbs, an orbit whose return is to a higher value of the bond than the one it left, drawn by an attractor that the system's own internal dynamics sustain. A forged cycle is the subtle and important case: it too may climb, but it climbs only because it is driven by an external input, a flow of energy or value across the boundary of the system, and is in the language of dynamics a dissipative structure, an order maintained at the expense of a throughput it does not itself generate. Cut the external input and the forged spiral collapses, whereas a genuine attractor persists on its own. This gives the forged-genuine distinction a precise dynamical signature: the genuine cycle is autonomous in the technical sense, sustained by its own vector field, while the forged cycle is non-autonomous, its apparent self-organization parasitic on a driving term from outside.

The limit of the dynamical formalization is the choice it cannot itself justify, and it is the choice of state variables. To write the dynamics one must first decide which quantities count as coordinates of the state, and that decision determines in advance what the model can see. Include only the bond, the recognition, and the consumption of the visible dyad, and the model will show a self-sustaining attractor, a genuine cycle, because the external streams that drive it have been omitted from the state space and reappear, if at all, only as unmodeled constants. Include the production labor, the unwaged self-fashioning, and the manufactured lack as state variables, and the same cycle reveals its driven, dissipative character. The dynamics does not choose its own state space; the modeler does, and the modeler's choice of what to include is precisely the drawing of the boundary whose stakes the political economy will make explicit. Dynamical systems thus offer an exact language for the distinction between autonomous and driven cycles and, in the same gesture, demonstrate that the distinction is only as honest as the state space is complete.

\subsection{Generative Grammar}

A second formalization treats \textit{amai} as a generative system, a set of production rules whose recursive application yields the unbounded variety of sweet expressions, much as a grammar yields the unbounded sentences of a language from a finite rule set. On this view the cuteness the second section analyzed is not a fixed inventory of charming things but a generative capacity, a small set of rules, diminutive forms, softened contours, marks of the unfinished and the appealing, that can be applied again and again to produce always new instances. The framework illuminates two features of the cycle with particular force. The first is its productivity: like a grammar, the sweet system can generate expressions it has never produced before, which is why the cycle never exhausts its repertoire and why consumption can always find a new object. The second is more pointed. The second section observed that the cute holds lack open, that it solicits a response it does not complete, and in grammatical terms this corresponds to a derivation that never reaches a terminal symbol, a production that always rewrites into something further to be rewritten. The endlessness of the cycle is, in this language, the non-termination of a derivation, and the framework explains formally why a desire mediated by the cute renews rather than concludes: the grammar has no stop rule.

The limit here is that a grammar specifies what can be generated, not whether what is generated is just. A generative system is normatively silent by construction; it characterizes well-formedness, the rules by which a sweet expression is correctly produced, and well-formedness is not justice. The same grammar generates the sweet expressions of a genuinely co-created intimacy and the sweet expressions of a capital-dictated repertoire, and nothing in the grammar distinguishes them, because the distinction lies not in the form of the derivations but in who authored the rules and at whose expense they are enacted, questions the formalism does not contain. Worse, and this is the framework's most useful self-disclosure, the choice of the rule set is itself prior to and outside the grammar. Whoever fixes the production rules fixes what sweetness can be, and the political economy's claim of real subsumption, that capital has reshaped the very grammar of \textit{amai}, is in this language the claim that capital has come to author the rule set itself. The grammar can show the generativity and the non-termination; it cannot, from within, say whose grammar it is.

\subsection{Field Theory}

The formalization closest to the series' own prior work treats the relation as a field and the cycle as transport around a closed path within it. Assign to the space of relational states a connection, a rule for how value is carried from one state to an adjacent one, and the holonomy of the cycle is the net transformation that value undergoes when carried once around the closed loop and back to its starting configuration. A vanishing holonomy is the flat case, mere circulation with no net change. A genuine positive holonomy is a real curvature of the relational field, an accumulation that the geometry itself produces, the formal counterpart of the spiral that climbs on its own. The forged holonomy, the paper's central diagnosis, acquires here its most exact definition: it is a phase that appears around the loop but is sourced from outside the region the loop encloses, a holonomy that is not the integral of an internal curvature but the trace of a flux entering across the boundary. This is the technical sense, developed in the author's work on geometric phase, in which the sweet cycle's apparent accumulation can be at once observable around the loop and absent as internal curvature, present as phase and forged as generation.

The limit of the field formalization is the one most worth stating, because the framework is the paper's own and the temptation to trust it as a master language is correspondingly strong. A field theory requires a choice of background and a choice of gauge. One must fix what counts as the space of states, what the connection is, and relative to what the phase is measured, and these choices are not delivered by the formalism; they are imposed on it. A holonomy is always a holonomy relative to a connection and a region, and to declare a cycle's holonomy positive or forged is already to have fixed the region the loop encloses, which is to say, once more, to have drawn the boundary. The field theory gives the forged holonomy a real and precise definition and simultaneously shows that the definition presupposes a boundary choice it cannot itself make. Even the paper's own most powerful formal instrument is, in this exact sense, not a view from nowhere.

\subsection{Category Theory}

Category theory shifts attention from objects to the structure-preserving maps between them, and it offers a language for the relations among the very frameworks this paper refuses to synthesize. One may regard each framework as a category, its objects the states of the cycle as that framework conceives them, its morphisms the transformations the framework recognizes, and one may then ask whether there exist functors, structure-preserving translations, between them. The categorial insight is twofold and cuts in both directions. On one side, a functor that faithfully carried the psychoanalytic category into the economic one, or the ethical into the dynamical, would be exactly the synthesis the paper denies is available; the absence of such faithful functors is the formal expression of incommensurability, the claim that there is no translation that preserves all the structure each framework sees. On the other side, category theory warns against a too-easy pluralism, for it also studies how partial translations, adjunctions and natural transformations, can relate frameworks without collapsing them, capturing the sense in which the frameworks are about the same cycle without being reducible to one another. The framework thus formalizes the paper's own methodological stance: many categories, no master functor, but structured relations among them all the same.

The limit is that category theory describes the form of relations among frameworks while remaining empty of their content, and in particular it cannot itself certify that any proposed functor is faithful or that any incommensurability is genuine rather than apparent. To assert that there is no master framework is, in categorial dress, to assert the non-existence of a certain functor, and that non-existence is not a theorem the formalism proves but a substantive claim about the frameworks that the formalism only lets us state cleanly. The danger, peculiar to this most abstract of tools, is to mistake having named the structure of incommensurability for having established it. Category theory gives the paper its sharpest way of saying what it means by incommensurable frameworks with structured relations; it does not relieve the paper of arguing, framework by framework, that the incommensurability is real.

\subsection{Game Theory}

Game theory models the sweet cycle as strategic interaction among agents each pursuing an outcome, and it captures with precision a dimension the other tools pass over, the dimension of incentive and equilibrium. Cast the firm, the indulged, and the indulgent as players with payoffs, and the cycle appears as a repeated game whose recurrence is explained by its being, in some range, an equilibrium: each party, given the others' strategies, does best by continuing. The framework explains formally why the cycle is stable without anyone intending its stability, why it persists as if by design though no one designed it, and it gives a precise content to the manufactured lack of the political economy, namely a deliberate alteration of the payoff structure by a party, the firm, that profits from keeping the others playing. The contemporary platform, analyzed earlier as a third party that accelerates the cycle, is in this language a player that designs the game others play, setting the payoffs so that the engagement it harvests is the equilibrium behavior of the rest. Game theory can even represent the forged character of the cycle as a negative externality: an equilibrium among the visible players that is sustained by imposing uncompensated costs on parties who are not at the table and whose payoffs the game, as drawn, does not include.

The limit, once again, is the boundary, and game theory makes it unusually visible because the boundary is literally the specification of who counts as a player. A game must fix its set of players, and whoever is left out of that set is, by construction, a party whose welfare the equilibrium analysis does not weigh, an externality at best. Draw the player set as the visible dyad and the firm, and the cycle is a benign equilibrium of mutual benefit. Admit the production worker and the bearer of manufactured anxiety as players, and the equilibrium is revealed as one sustained by externalized costs, a solution that is stable only because the losers have no seat. The choice of the player set is the drawing of the boundary in the most literal form the paper will encounter, and game theory, in requiring that choice before it can compute anything, demonstrates yet again that the formal verdict is hostage to a prior and non-formal decision about who is admitted to the reckoning.

\subsection{Information Theory}

Information theory offers a last formalization, treating the sweet cycle as a channel along which a signal, the confirmation of being wanted, is transmitted from one party to another through the medium of the object. The framework lights up the second section's analysis of the demand for love from an unexpected angle. What the indulged party seeks is, in these terms, a reduction of uncertainty about the other's desire, a signal that lowers the entropy of the question \textit{am I still wanted}, and the cute object functions as a carrier of that signal. The endlessness of the cycle becomes legible as a property of a channel in which the uncertainty is never fully resolved, in which each message reduces the doubt only partially and the doubt regenerates, so that the channel must carry the signal again. The framework also gives a precise way to state what capital does: it inserts itself as the owner of the channel, charging for the transmission of a signal that once passed freely, and it can even degrade the channel deliberately, manufacturing noise, the produced sense of insufficiency, so that more messages, more purchases, are required to achieve the same reduction of uncertainty. The manufactured lack is, in this language, manufactured entropy, noise introduced into the channel to increase the volume of paid transmission needed to overcome it.

The limit of the information-theoretic picture is that it measures quantity of information and says nothing of its worth. A channel's capacity, its noise, its rate are all definable without any reference to whether the signal transmitted is true, whether the reassurance carried is honest, or whether the channel is justly owned. Information theory can model the reduction of uncertainty about being wanted and remain wholly indifferent to whether one is in fact wanted or merely told so, and it can model the manufactured noise without judging the manufacture. As with every tool in this section, the formalism requires a prior choice, here of what counts as signal and what as noise, and that choice is not innocent: to classify the produced sense of insufficiency as noise is already a normative judgment that the formalism borrows from outside itself, for from a purely information-theoretic standpoint the manufactured anxiety is simply more signal in a different channel, the channel of the market. The tool measures the flow and is silent on its justice and its truth.

\subsection{The Reflexive Lesson: Formalization Is a Choice of Boundary}

A pattern has repeated across all six tools, and it is the reason this section belongs in a paper about justice rather than in a technical appendix. Each formal instrument captured something the substantive frameworks could not state with comparable precision: the dynamical signature of a driven versus an autonomous cycle, the non-termination of the cute derivation, the exact definition of a forged holonomy, the structural meaning of incommensurability, the equilibrium logic of stability and externality, the entropic account of manufactured lack. And each, at the point of its greatest precision, disclosed that its precision rested on a prior choice that the formalism itself could not make: the choice of state variables, of rule set, of background and gauge, of which functors to admit, of who counts as a player, of what counts as signal. In every case the choice was, in substance, the drawing of a boundary, a decision about what the representation includes and what it consigns to the outside.

This is the reflexive lesson, and it closes the loop with the rest of the paper. The political economy will argue that the judgment of the cycle turns on where the boundary is drawn and who is allowed to fall outside it. The formal tools, far from supplying a boundary-free vantage from which that judgment could be made objectively, are revealed to be themselves so many ways of drawing the boundary, each making its cut and computing its verdict relative to the cut it made. There is no formalization that is not a framing, and no framing that does not include and exclude. This does not make the formal tools worthless; it makes them honest, once their boundary choices are made explicit, and it forbids the particular error of treating any of them as the master language that would adjudicate from above what the substantive frameworks left open. The most rigorous instruments in the paper turn out to confirm, rather than escape, its central methodological claim: that the goodness of a cycle that looks good is decided differently depending on where one draws the line, and that drawing the line is an act for which no formalism can be held responsible in place of the one who draws it.

\section{Political Economy: How Capital Produces, Reproduces, and Extracts from Amai}

The frameworks gathered so far have established that the satisfactions of the sweet cycle are structurally intelligible, neurally real, and, within the relation, capable of being a genuine practice of care. None of them has been able to say whether the cycle is just, and two of them, psychoanalysis and neuroscience, have explicitly declined the question as outside their competence. This section takes up the question they decline. Its thesis is that the sweet cycle, viewed at the scale at which value is produced and reproduced, is not a closed loop that generates its own surplus but an open system that imports its surplus from outside while presenting itself as closed, and that this importation, hidden by the very framing the cycle invites, is extraction in the precise sense the series has developed.

The analysis proceeds in three layers, each deeper than the last. The first is the classical analysis of surplus and social reproduction, which locates where the cycle's value is drawn from. The second is a philosophical analysis in the Marxian tradition, which asks what the cycle does to the persons caught in it, through the categories of alienation, fetishism, subsumption, reification, and ideology. The third is an analysis in the idiom of contemporary political economy, which asks what becomes of the cycle once it is organized by emotional capitalism, by platforms and the attention economy, by the aesthetic economy, and by the rentier logic that increasingly governs accumulation. The three layers are not rivals; each sees what the others cannot, and together they show that the diagnosis of a forged holonomy is not the verdict of a single school but a conclusion reached, by different routes, from several.

\subsection{The Mechanism, Stated Plainly}

Before the layers, the central mechanism deserves to be stated in one place and at its full sharpness, because it is the spine on which everything in this section hangs and it is more precise, and more damning, than the familiar charge that capital manufactures false desires. That familiar charge is weak, because it presupposes a contrast with some catalogue of true desires by which the false ones could be exposed, and no such catalogue survives scrutiny: the wish for a small bright thing is no more counterfeit than the wish for bread. The mechanism this paper describes does not depend on any such contrast and is the stronger for it. What capital has done is not to fabricate a desire that was not there. It is to discover a structural truth about desire as such, the truth the second section reconstructed, that a living relation is sustained not by satisfaction but by a lack kept in operation, with the cute serving as the cultural mediation through which that lack is approached without being filled, and then to occupy the position that truth defines. The move has three beats. First, capital inserts its commodity into the place where the mediation of lack already does its work, so that the recognition which formerly flowed through a glance or a word must now flow through a purchase. Second, because the mediation of lack is by its nature unclosable, the openness that makes the relation alive, the purchasing it now carries is likewise unclosable, and the cycle renews without end not because anyone is addicted but because the structure of desire was never going to terminate. Capital does not need to manufacture endlessness; it inherits it from the form it has occupied. Third, and most consequentially, capital does not merely occupy the lack but enlarges its production, reproducing on a widening scale the sense of not yet being sweet enough, because the manufactured insufficiency is the fuel the cycle burns. The very feature that the reconstruction identified as the condition of a living relation, the persistence of lack, is thus taken over, industrialized, and made to yield a surplus.

The paradox that the prologue felt at the magazine cover is now explicable, and it is the heart of the matter. The cycle looks just precisely because it is built upon a real and nourishing structure rather than upon a deception. The satisfactions are genuine, as the neuroscience confirmed; the care is genuine, as the ethics of care insisted; the mediation of lack really is keeping the relation alive. The cover's claim that sweetness is justice is therefore not a simple lie that a sharper eye would see through, because within the cycle's own boundary it is, in part, true. This is exactly why the critique cannot proceed by unmasking a falsehood. It must proceed by widening the boundary, by following the value across the edge the cycle prefers we not cross, and asking whether the surplus that circulates so warmly within was generated within or drawn in from a labor and a lack kept just out of frame. The genius of the arrangement, and the reason it can wear the name of justice without obvious absurdity, is that it has wrapped a genuine intimacy around an extraction so that the two present a single sweet face. The work of the analysis that follows is to show that the face has a hidden side, and to say precisely whose labor and whose manufactured want are kept there.

\subsection{The Classical Layer: Surplus, Social Reproduction, and the Forged Holonomy}

The classical analysis proceeds in four movements. The first movement concerns the commodity mediation of a demand that is not, in itself, a commodity. What the indulged party seeks, the earlier sections established, is the confirmation of being wanted, a demand for love that operates in the register of the symbolic and cannot be directly bought or sold. Capital's operation is to insert a commodity into the path of that demand. The cosmetic, the accessory, the small sweet thing becomes the necessary mediation through which the confirmation is sought and given. A circuit that in itself runs on recognition is thereby rewritten as a circuit that runs on purchase. In the series' shorthand, a demand $D$ is routed through a commodity $C$, and what had been an exchange of recognition becomes a movement of money through goods and back to money enlarged, the form the series has written as $M$-$C$-$M'$. The crucial point is that this rewriting is not a neutral convenience. It makes the continuation of the recognition depend on the continuation of the purchasing, and it is from that dependence that everything else follows.

The second movement asks where the surplus comes from, and it is here that the analysis must be most concrete, because the difference between a genuine and a forged cycle is nothing other than the answer to this question. If the surplus that the cycle accumulates, the deepened bond and the firm's enlarged $M'$, were generated entirely within the loop by the participants' own exchange, the cycle would be genuine. It is not generated entirely within the loop. At least three streams of value flow into the cycle from across its edge, and all three are systematically rendered invisible by the framing that lets the cycle look closed. The first is the labor of the production chain, the making of the cosmetics and the garments, frequently performed far from the site of consumption, frequently by women, under conditions the sweet cycle never has to see. The second, and the one this paper most wants to make visible, is the reproductive labor that the indulged party performs upon themselves, the work of producing and maintaining the very image of being worth indulging, the time, money, attention, and self-management that go into being sweet. This labor is unwaged, and more than unwaged, it is naturalized as self-care and even as self-love, so that its character as labor disappears precisely in being praised. The third stream is the manufactured lack itself. Capital does not find a fixed quantity of desire and serve it; it continually produces the sense of not yet being sweet enough, not yet indulgent enough, not yet worthy, because that produced insufficiency is the engine that keeps the purchasing in motion. Here the analysis rejoins the author's earlier work on lack: capital has discovered the structural truth that a relation runs on a lack kept in operation, and it has industrialized the production of that lack.

The second of these three streams repays a closer look, because it is the one most thoroughly hidden and the one the present series is best positioned to make visible. The work of producing oneself as worth indulging, the maintenance of the sweet self, belongs to the broad category that feminist political economy has taught us to call social reproduction: the vast, largely unwaged, largely feminized labor of producing and maintaining the persons who then appear in the waged economy as if already made. Nancy Fraser's formulation is useful here, that capital depends upon conditions of social reproduction which it does not pay for and indeed tends to erode, consuming the very care and relational capacity that it presupposes. The sweet cycle is a precise instance of this larger structure. The capacity to be sweet, to present an appealing and indulgable self, does not arise from nowhere; it is produced and continually reproduced by labor, and that labor is performed off the books, by the indulged party upon themselves, in hours that no wage counts and that the culture reclassifies as leisure, pleasure, or self-love. The genius of the arrangement, from capital's standpoint, is that the reproductive labor on which the cycle depends is not merely unpaid but is experienced by the one performing it as a gift to herself, so that the extraction is felt as indulgence and the cost as a treat. There is no clearer case of what it means for an extraction to be hidden by the very framing that the commodity supplies, and it is no accident that the framing the Ginza cover supplies is, precisely, that this labor is justice.

The third stream, the manufactured lack, completes the picture and connects the political economy back to the psychoanalysis. The earlier sections established that desire runs on a lack kept open, and insisted that this openness is in itself the condition of a living relation rather than a defect. The political economy adds the decisive observation that this openness is a resource which capital has learned to manufacture and to enlarge. The produced sense of insufficiency, of not yet being sweet enough, is not a byproduct of the cycle but its fuel, deliberately generated because a satisfied subject buys nothing. What capital exploits is therefore not a contingent insecurity that better marketing happened to create, but the structural openness of desire itself, harnessed and amplified into a permanent and expanding demand. This is why the critique cannot end by recommending that people simply feel more secure, as though the manufactured lack were a correctable error rather than the operating principle of the arrangement. The lack is doing exactly what it is built to do.

The third movement names the result, and the name is one the series has prepared. A cycle whose inner surplus is sustained by streams of value drawn from across its edge, while it presents itself as a closed loop generating its own surplus, is a forged holonomy. The positive accumulation that the strongest version of the good cycle correctly perceived is real as an appearance and false as an accounting, because the books were balanced only by leaving the external streams off them. Let the boundary be drawn where the strongest version drew it, around the firm, the indulged, and the indulgent, and the surplus reads as positive, value generated and returned to its generators. Let the boundary be drawn wide enough to include the production worker, the unwaged reproductive labor of self-fashioning, and the manufactured anxiety, and the sign of the surplus is no longer obviously positive, because some of what the inner parties enjoy is not generated by their exchange but transferred into it from those who bear its costs and are counted in none of its benefits. In the series' terms, the difference between a genuine and a forged cycle is the difference between a closed system that accumulates a real internal holonomy and an open system that fakes one by extraction, and the sweet cycle, at the scale of its production and reproduction, exhibits the second form. This is the same diagnosis the series has applied to the general formula of capital, now made specific to the intimate register, where it is harder to see and therefore more worth seeing.

It may help to state the diagnosis in the more quantitative idiom the series uses elsewhere, while being honest about what the idiom can and cannot deliver in a case like this. The series writes the net value a cycle accumulates over a traversal as the difference between what it generates and what it consumes, a quantity whose sign distinguishes a genuinely accumulating cycle from a merely circulating or a depleting one. The crucial subtlety, and the whole point of the forged-holonomy diagnosis, is that this quantity is not boundary-independent: its sign depends on where the edge of the cycle is drawn, because the drawing of the edge determines which generative contributions and which consumptive costs are counted. Drawn narrowly, around the visible participants, the sweet cycle shows a clearly positive net, much is generated and little consumed, because the largest costs, the production labor, the unwaged self-fashioning, the manufactured anxiety, fall outside the boundary and are not entered on the ledger. Drawn widely, to include those costs, the net is no longer reliably positive, because much of what appeared as internally generated value is revealed as value transferred in from across the edge. I do not pretend that this quantity can be measured in any precise empirical sense for a phenomenon as diffuse as \textit{amai}; the point of writing it is not to compute it but to make visible that the appearance of a positive holonomy is an artifact of a boundary choice, and that the boundary choice is precisely what the cycle's self-presentation is designed to naturalize. The forged cycle is, in this exact sense, a cycle whose positive sign survives only so long as the accounting refuses to widen.

The fourth movement observes that the cycle does not merely repeat but expands, and that its expansion is its most telling feature. Each turn of the sweet cycle does more than strengthen the bond; it trains the participants to bind recognition more tightly to consumption, so that the next confirmation of being wanted requires a little more than the last. The lack that capital produces is not produced once but reproduced on an enlarging scale, and the recognition that was once available through a small mediation comes to require a larger one. This is the intimate form of accumulation, $M'$ greater than $M$ carried not by the firm's ledger alone but by the deepening dependence of recognition upon purchase. And it is exactly here that the forged cycle reveals its difference from a genuinely generative one. A genuinely good cycle, in the sense the series has defended, would accumulate its positive holonomy internally and would not need to draw an ever-larger tribute from outside to do so; its spiral would climb on its own generated value. The sweet cycle climbs, when it climbs, by reaching further across its edge each time. The expansion that looks from inside like a deepening love is, at the scale of the system, a widening extraction. To say this is not to say that no one inside is really in love. It is to say that the love and the extraction are, in the forged cycle, running on the same track, and that the work of the practical sections later in this paper will be to ask whether they can be pulled apart.

\subsection{The Philosophical Layer: Alienation, Fetishism, Subsumption, Reification, Ideology}

The classical layer locates where the value is drawn from. It does not yet say what the cycle does to the persons who live inside it, and for that the analysis must move from Marx's economics to his philosophy, and to the tradition that extended it. Five categories are needed, each illuminating a different facet of what the sweet cycle does to subjectivity itself. They are not loose metaphors borrowed for color; they are precise concepts, and the sweet cycle turns out to be an unusually clean instance of each.

The first is alienation, in the sense of the early Marx of the 1844 manuscripts, where the worker's own product confronts her as an alien power standing over against her and governing her. The sweet cycle produces an intimate variant of exactly this. What the indulged party produces, through the unwaged labor of self-fashioning, is an image of the sweet self, and that image, once produced, confronts her as something she must continually live up to, an alien standard that governs her rather than a free expression that she governs. She does not possess the sweet self; the sweet self possesses her, dictating what must be bought, maintained, performed, and renewed. The product of her labor has become a power over her. This is alienation in the strict sense, and it connects directly to the imaginary capture diagnosed in the second section: the image that the imaginary fixes is precisely the alienated product, the self externalized into a picture that then commands its original. The classical concept and the psychoanalytic one name, from two directions, a single structure.

The second is commodity fetishism, in the sense of the first volume of \textit{Capital}, where a definite social relation between people assumes the fantastic form of a relation between things. The sweet cycle is fetishistic to its core. The relation between the indulged and the indulgent, and behind them the relation between the consumer and the distant workers who made the objects, appears in the cycle as a relation between a person and her possessions, between a self and the bright things that make it sweet. The social relations, the labor of the production chain, the reproductive labor of self-fashioning, vanish into the object and reappear as a property of the object itself, its sweetness, its cuteness, its power to charm. The Ginza cover is the perfect fetishistic formula. To say that sweetness is justice is to attribute to a commodity, to the quality of the sweet thing, a predicate, justice, that properly belongs to a social relation among persons. The fetish does not merely hide the social relation; it relocates its qualities into the thing, so that one seeks in the object what can only be found in the relation. This is why the cycle can feel like the pursuit of justice while being the consumption of goods: the fetish has lodged the social predicate in the commodity.

The third is the distinction between formal and real subsumption, from the unpublished sixth chapter of \textit{Capital}. Capital formally subsumes a labor process when it takes over a process that already existed and sets it to producing surplus without changing its inner form; it really subsumes that process when it reshapes the process itself to its own requirements, remaking the activity from within. The sweet cycle is a case of real subsumption reaching into the intimate sphere. Capital did not merely find an existing practice of \textit{amai} and attach commodities to it, which would be formal subsumption. It has reshaped \textit{amai} itself, remade what it is to be sweet, to be indulged, to charm and be charmed, around the commodity form, so that the very grammar of intimate reliance now bears the impress of what must be bought to enact it. The standard of sweetness, the tempo of its renewal, the forms it may take, are increasingly dictated by the commodity system rather than by the relation. Intimacy has not merely been monetized at its edges; it has been reorganized from within. This is the deepest sense in which the cycle is not a natural phenomenon that capital exploits but a phenomenon that capital has, in part, produced.

The fourth is reification, in Lukács's extension of fetishism, where the commodity form so saturates consciousness that relations, processes, and qualities that are living and dynamic come to be apprehended as fixed things. The sweet cycle reifies intimacy. \textit{Amai} is a process, a continual mutual generation of recognition, a dynamic that lives only in its enactment. Under the commodity form it is frozen into a set of possessable things, the look, the items, the image, that can be acquired, displayed, and exchanged. The living dialectic of a relation that the second section described, the open question continually reposed, is congealed into an inventory. What was an activity becomes a stock. And because the reified form is what the market can sell, the market continually pulls the relation toward its reified version, toward the having of sweetness rather than the doing of it. Reification is the bridge between the fetishism of the object and the alienation of the self: the relation is thingified, the self is thingified into its image, and both then circulate as commodities.

The fifth concerns ideology, and here the contemporary form of the critique, associated with Žižek, is sharper than the classical one. The classical theory of ideology held that the participants are deceived, that they do not know the real conditions of their activity and would withdraw if they did. But the sweet cycle is not sustained by ignorance. Many who participate know perfectly well that they are inside a consumerist apparatus, can articulate the critique, and continue all the same. This is the cynical structure Žižek described: they know very well what they are doing, but still they do it. Ideology here operates not at the level of knowledge but at the level of practice, in the fetish that is enacted regardless of what is known. The implication for this paper is important and is the reason the critique cannot rest on consciousness-raising alone. Telling the participants the truth about the cycle does not dissolve the cycle, because they were not held in it by a falsehood in the first place. They are held by a practice, by the lived dependence of recognition upon purchase, and a practice is not undone by being seen through. This is why the paper's constructive sections must speak of practices and not merely of awareness: only a changed practice answers a fetish enacted in spite of knowledge.

Taken together these five categories deliver a verdict the classical layer could not. The classical layer showed that the cycle extracts. The philosophical layer shows that it deforms: it alienates the self into an image that rules it, relocates the qualities of social relations into commodities, reshapes intimacy from within to the measure of the market, freezes a living relation into possessable things, and secures participation through a fetish that survives full knowledge. The extraction and the deformation are two aspects of one process, and neither alone is the whole of what is wrong.

\subsection{The Contemporary Layer: Emotional, Platform, Aesthetic, and Rentier Capitalism}

The categories so far are powerful but were forged for an industrial capitalism of factories and waged labor. The sweet cycle lives in a different economy, and four contemporary developments in political economy are needed to see how the old mechanism operates under new conditions. Each names a form of accumulation that the classical analysis did not anticipate and that the sweet cycle exemplifies.

The first is emotional capitalism, in Eva Illouz's sense, the historical process by which emotional life and economic life have interpenetrated, so that emotions are made into objects of calculation, management, and exchange, while economic relations are saturated with the language of feeling. The sweet cycle is emotional capitalism in miniature. The demand for love that the second section placed at its center is, under emotional capitalism, recoded as something to be managed by consumption, optimized, invested in, and audited for its returns. The intimate question, am I still wanted, becomes a question with a purchasable answer and therefore a budget, a strategy, a portfolio of sweet things assembled to secure an emotional yield. What Illouz lets us see is that this is not a corruption of an otherwise pure emotional sphere by an external economic logic; the two have grown together, so that the very way the participants experience their longing is already shaped by the market that offers to satisfy it. The demand for love arrives pre-formatted for consumption.

The second is the platform and attention economy, and here the analysis must update its picture of where the sweet cycle now runs. The Ginza magazine cover is in a sense already an old medium. The contemporary sweet cycle runs substantially on platforms, where the cute self is performed for an audience and the performance is captured as data. Two extractions compound here. In the terms of Shoshana Zuboff's analysis of surveillance capitalism, the performance of the sweet self generates a behavioral surplus, a stream of data about attention, desire, and response that the platform appropriates and turns to account, so that the unwaged labor of self-fashioning now produces, in addition to the image, a raw material that is extracted at a second remove. And in the terms of the attention economy, attention itself is the scarce good, and the cute, which the second section showed to be a form that solicits response, is an unusually efficient instrument for capturing it. The platform is therefore not a neutral venue for the sweet cycle but an active party to it, engineering the cycle's renewal because the cycle produces the engagement the platform sells. The intimate dyad of the original analysis is now embedded in a triad whose third member, the platform, profits from the dyad's every turn and has every incentive to accelerate it.

The third is aesthetic capitalism, in Gernot Böhme's sense, the stage of accumulation in which the production of use values gives way increasingly to the production of staging values, the value a commodity has in letting its owner stage an appearance, an atmosphere, a self. The sweet things of the cycle are aesthetic commodities par excellence. What they sell is not a use but a staging, the power to appear as the sweet self, to compose an atmosphere of indulgability. Böhme lets us name precisely what is consumed: not the cosmetic as pigment but the cosmetic as the means of staging a self worth charming. This connects the contemporary layer back to the reification of the philosophical layer and the imaginary capture of the psychoanalytic one, for the staging value is exactly the price of the alienated image. Aesthetic capitalism is the economic form whose product is the imaginary itself, sold by the unit.

The fourth is the rentier logic that, on a growing body of analysis, increasingly characterizes contemporary accumulation, the shift from profit earned by producing value to rent extracted by controlling an asset that others must pay to use. The relevant asset in the sweet cycle is the brand, and behind the brand the whole apparatus of signs, the cuteness, the names, the aesthetic codes, that the consumer must pay to deploy. The premium paid for the branded sweet thing over its cost of production is, in this light, not a return to value the firm created but a rent on a sign the firm controls, a charge for access to a vocabulary of sweetness that has been enclosed. The distinction, central to the work of writers such as Mazzucato on value creation versus value extraction, lets us sharpen the forged-holonomy diagnosis one final degree. To the extent the firm's enlarged $M'$ is rent rather than profit, a charge for access to an enclosed sign rather than a return to generated value, the cycle is forged not only because it draws on external labor but because its central commercial party is extracting rather than creating, collecting a toll at a gate it has erected across the path of the demand for love. The enclosure of sweetness is the rentier moment of the sweet cycle.

These four contemporary forms do not replace the classical and philosophical analyses; they show those analyses operating under present conditions and intensified by them. Emotional capitalism preformats the demand, the platform extracts a second surplus from its performance and accelerates its renewal, aesthetic capitalism sells the staging of the alienated image, and rentier logic collects a toll on the enclosed signs through which all of this must pass. The sweet cycle of the Ginza cover, read through these four, is revealed as a node in an accumulation regime far larger and more extractive than the intimate dyad it appears, from inside, to be.

\subsection{A Caution on What the Diagnosis Does and Does Not Claim}

I close this section with a caution the series requires of itself and which must not be omitted. To diagnose the sweet cycle as a forged holonomy is a structural claim, not a verdict on any person within it, and it does not by itself establish that the satisfactions are false or the care insincere. The earlier frameworks have already shown that they are neither. What the structural claim establishes is narrower and harder: that the goodness perceived from inside is, in the cycle as capital currently organizes it, underwritten from outside, and that an honest reckoning cannot stop at the satisfied faces in the frame. Whether a sweet cycle could be organized otherwise, so that its surplus were genuinely internal and no one across its edge were reduced to fuel, is not a question political economy can answer by critique alone. It is a question for practice, and the paper will reach it.

\section{Feminism: Extraction, or a Stigmatized Capacity for Relation?}

No framework divides against itself before the sweet cycle as feminism does, and the division is not a weakness to be resolved but the most instructive thing feminism has to offer this inquiry. The cycle presents feminism with a phenomenon that its two deepest commitments read in opposite directions, and holding both readings at full strength, rather than collapsing one into the other, is the discipline this section requires.

The first reading is the critical one, and it converges with the political economy just developed. \textit{Amai} is not distributed evenly across persons; it falls along a gendered line. It is overwhelmingly women who are addressed by the Ginza cover, women who are enjoined to be sweet, cute, worth indulging, and women who perform the unwaged reproductive labor of producing that sweetness upon themselves. The critical tradition has a precise vocabulary for what is happening. The management of feeling and appearance to meet an external standard is labor, emotional labor in Arlie Hochschild's sense, performed continuously and counted nowhere. The standard of cuteness to which the labor is addressed is a disciplined femininity, a norm that produces the very subjects who then experience meeting it as their own desire. And the manufactured anxiety that drives the purchasing, the produced sense of never being sweet enough, falls disproportionately on women as a continuous tax on self-worth. Read this way, the proposition that sweetness is justice is a near-perfect specimen of ideology: it confers the name of a virtue on a gendered extraction, so that the labor of being sweet is experienced not as a cost imposed but as a justice embraced. This reading is powerful and, within its scope, correct.

The critical reading has a further resource that sharpens it beyond the charge of mere unfairness, and it is worth drawing out because it connects the gendered analysis to the imaginary capture diagnosed in the second section. The femininity that the cycle rewards is not a pre-existing trait that the cycle happens to remunerate; it is, in significant part, produced by the repeated performance the cycle demands. The norm precedes the subject and shapes the desire through which the subject then experiences the norm as her own. This is why the appeal to the participant's felt willingness cannot settle the question in the cycle's favor. That a woman wants to be sweet, enjoys the small bright things, experiences the labor of self-fashioning as pleasure, is not evidence that the cycle is innocent, because the wanting, the enjoyment, and the experience of pleasure are themselves among the cycle's products. The critical reading thus reaches the same wall the neuroscience reached, from the other side: subjective endorsement, however sincere, cannot certify justice, because the endorsement is part of what the arrangement manufactures. Where the neuroscientist said that real happiness does not entail a just cycle, the feminist critic adds that the happiness may be the cycle's most effective instrument, the very means by which a gendered extraction secures the consent of those from whom it extracts.

This is the critical reading at full strength, and the temptation at this point is to take it as the conclusion. The section's discipline is to refuse that temptation, not by weakening the critique, which is sound as far as it reaches, but by insisting that it does not reach far enough, and that a second reading, equally feminist, sees something the first is structurally liable to miss.

The second reading refuses to let the first have the last word, and it does so on feminist grounds, not against them. To read \textit{amai} solely as women being disciplined is to perform a familiar and suspect gesture, the gesture of denying women's agency in the name of protecting them, of treating what women do as something merely done to them. The ethics of care prepared the counter-reading: dependence, softness, the solicitation of indulgence are not in themselves marks of subordination but capacities for relation that a culture organized around masculine-coded autonomy has taught us to despise. To presume upon another's care, and to do so deliberately, skillfully, even strategically, can be an exercise of power within a relation rather than a surrender of it. The one who is sweet is not always the one who is weak. There is a tradition of feminist thought, and there is a complexity specific to the Japanese context that resists any simple importation of an Anglophone narrative of disciplined femininity, that insists \textit{amai} can be a competence women wield, a register in which they act, and not only a mold into which they are pressed. To deny this is to repeat, in critical dress, the same devaluation of the relational that the critique elsewhere opposes.

The point about agency must be made carefully, because it is easily heard as a capitulation, as though to grant that women act within the cycle were to grant that the cycle is fine. It is not that. The second reading does not deny the extraction the first reading documents; it denies that documenting the extraction exhausts the description of what women are doing. A woman who deploys sweetness can be, at one and the same time, the object of a gendered extraction and the agent of a relational strategy she has mastered, and to see only the first is to reproduce the condescension that has always accompanied the disparagement of feminine-coded skill. The history of devaluing care, receptivity, and the management of relation as not-quite-work, not-quite-skill, not-quite-agency is itself a feminist target, and a critique of \textit{amai} that treats its practitioners purely as the patients of an imposed norm has, despite its intentions, joined the side that cannot see relational competence as competence. The second reading insists that the same act can carry both descriptions and that the refusal to let either silence the other is what fidelity to the phenomenon requires.

The Japanese context sharpens rather than dissolves this, and a brief note guards against a too-easy transposition. The Anglophone narrative of disciplined femininity arrives with its own assumptions about autonomy, individuality, and the suspicion of dependence, assumptions that are not culturally neutral and that the concept of \cn{甘え} was in part articulated by Doi to contest. To read \textit{amai} through a framework that treats all dependence as presumptively suspect is to risk importing precisely the autonomy-fetish that the ethics of care and Doi's analysis both call into question, and thereby to misdescribe as subordination what may, in its own cultural grammar, be a recognized and reciprocal mode of relation. This is not an argument that the gendered extraction is absent in Japan, it plainly is not absent, but a caution that the line between extraction and competence cannot be drawn by an imported template and must be drawn, if it can be drawn at all, with attention to the specific relational grammar in which \textit{amai} actually operates. The cultural particularity that a later section addresses in general terms returns here in a pointed form: even the feminist judgment is boundary-dependent, and the boundary in question is partly cultural.

The two readings do not cancel. They locate a real question that neither alone can answer, and stating that question precisely is this section's contribution to the judgment that follows. The question is this: when is the gendered asymmetry of \textit{amai} an extraction, the channeling of women's unwaged reproductive labor into a cycle that profits elsewhere, and when is it a stigmatized capacity for relation, a genuine competence that a culture of autonomy has wrongly taught us to read as weakness? These are not two names for one thing. They are two different things that can wear the same appearance, and the difference between them is, once again, the difference between a cycle that returns its value to those who generate it and one that does not. Feminism, divided against itself, delivers the inquiry's sharpest formulation of its own central question and, characteristically, refuses to pretend it can settle it from theory alone.

\section{The Test of Judgment: Eudaimonic? Just? Genuine or Forged?}

The frameworks have spoken and have not agreed. This section does not overrule them. It gathers the questions they raised into three criteria that can be stated and applied here, in this paper, without waiting on the fuller development they receive elsewhere.\footnote{The three criteria are given here in a deliberately minimal, self-standing form sufficient for the present case. Their full articulation, including the argument that sublimation rather than exhaustion is what a positive phase measures, and the conditions under which the criteria can be jointly applied, belongs to a companion paper in this series on the practice of flourishing, in preparation. Nothing in the present application depends on that fuller account.} The criteria are not a master framework that adjudicates the others; they are a way of asking, of any concrete cycle, three questions whose answers the frameworks have equipped us to seek.

The first criterion is the eudaimonic one, and it can be stated in the classical register the series has used before. A cycle is eudaimonic to the degree that it lets each participant's own singular dynamic unfold, in the Aristotelian sense in which flourishing is the activity of a life according to what is distinctively its own, rather than fixing each participant into a flattering but static image. Applied to the sweet cycle, the question is the one armed in the second section. Does the cycle open the slow, generative work of two singular dynamics meeting and unfolding, or does it substitute for that work the consumption of an image of being indulged, capturing the living relation in the imaginary and selling the picture back in place of the thing? The frameworks give a divided answer, which is the honest one. To the extent the objects mediate a real address between persons, the cycle can be eudaimonic. To the extent the image of the sweet self comes to be consumed in place of the unfolding it pictures, the cycle forecloses flourishing while appearing to deliver it.

The second criterion is justice, and it can be stated as a single operational test that the political economy section has already put to work. A cycle is just only if no one across its edge is reduced to mere fuel, contributing the value the cycle enjoys while counted in none of its benefits. This is not a demand that a cycle have no outside; every cycle has an outside. It is a demand that the outside not be a reservoir of unreturned value, that those whose labor and whose manufactured lack sustain the inner surplus not be permanently external to the distribution of what they sustain. By this test the sweet cycle, as capital currently organizes it, does not pass, because the production worker, the unwaged self-fashioning, and the manufactured anxiety are precisely such unreturned external streams.

The third criterion subsumes the first two into the structural distinction the series turns on, between a genuine and a forged accumulation. A cycle is genuine when the value it accumulates is generated within it and the spiral climbs on its own internally generated surplus; it is forged when the inner accumulation is sustained by importation from across the edge while the cycle presents itself as closed. The sign of the difference is whether the internal surplus survives an honest widening of the boundary. Widen the boundary of the sweet cycle and, as the previous sections showed, the surplus does not survive intact; it was in part transferred, not generated. The cycle is, in its current organization, forged.

It is worth being explicit about how the three criteria relate, since they are not simply three items on a checklist. The eudaimonic criterion concerns the quality of the relation for those within it, whether it opens or forecloses the unfolding of singular lives. The justice criterion concerns the relation's outside, whether it leaves anyone across its edge as unreturned fuel. The genuine-versus-forged criterion is the structural one that, in a sense, contains the other two: a cycle is forged exactly when its inner goods, including whatever eudaimonic and ethical goods it really does deliver, are sustained by an unjust importation from outside. This is why the three answers can diverge without contradiction. A cycle might be eudaimonic for its participants and yet forged, because the flourishing it affords them is paid for by an extraction they do not see; this is, the paper has argued, the actual situation of \textit{amai} as capital now organizes it. The criteria are therefore not redundant and not reducible to one another; they measure different things, and a complete judgment requires all three precisely because a cycle can pass one while failing another. The refusal to collapse them into a single score is not indecision but accuracy, a recognition that the goodness of a cycle is irreducibly plural in its dimensions.

These three answers do not combine into a single verdict, and the refusal to combine them is not evasion. The cycle is, at once, capable of being eudaimonic in the meeting of persons, unjust in the external streams it leaves unreturned, and forged in the structure of its accumulation. The right conclusion is not that one of these is the truth and the others appearance. It is that the goodness of a cycle that looks good is decided differently at different scales, and that the decision turns, in every case, on where the boundary is drawn and who is allowed to fall outside it. That observation does not end the inquiry. It hands it to practice, which is the only place where a boundary can actually be redrawn.

\section{Cultural Particularity and the Question of the Universal}

An honesty the inquiry owes itself concerns the status of its central example. \textit{Amai} is not a culturally neutral specimen of intimate consumption. It is a structure that Doi raised to theoretical prominence as something characteristically, though he did not claim exclusively, Japanese, and the Ginza cover that occasioned this paper is a particular artifact of a particular consumer culture. The question this raises is whether the judgments reached above are about \textit{amai}, this culturally specific structure, or about intimate consumption cycles as such, and the paper should not let the convenience of its example obscure the difference. Before drawing the methodological conclusion, the cultural specificity itself deserves substantive attention, for the inquiry has leaned heavily on \textit{amai} and \textit{kawaii} without yet asking why these forms are so culturally pronounced where they are.

\subsection{The Cultural Specificity of Amae}

That dependence might be regarded not as a failing but as a recognized and even valued mode of relation is the central claim of Takeo Doi's \textit{The Anatomy of Dependence}, and the claim is explicitly comparative. Doi developed the concept of \textit{amae} in part by contrast with the assumptions about autonomy and independence he encountered in Western, and especially American, clinical and social settings, arguing that Japanese social life affords \textit{amae} a structural centrality and a positive valuation that the autonomy-centered cultures tend to withhold. One need not accept every element of Doi's account, and much in it has been contested, including the risk of essentializing a complex society into a single relational key and the nihonjinron tendency, the genre of writing on Japanese uniqueness, into which his work was sometimes absorbed. The durable point that survives the criticism is narrower and sufficient for this paper: that whether dependence reads as weakness or as competence is not a cultural universal but varies with the relational vocabulary a society makes available, and that Japanese culture has historically made available a vocabulary in which licensed reliance is legible as a bond rather than a deficiency. This is precisely the cultural condition under which a magazine can place sweetness and justice in apposition and be understood, rather than dismissed, by its readers.

\subsection{Kawaii as a Historically Specific Formation}

The aesthetic of \textit{kawaii}, the cute, on which the sweet cycle so heavily draws, has itself been the object of serious cultural study, and that study delivers a result this paper needs: \textit{kawaii} is not a timeless taste but a historically specific formation, which is exactly what one would expect of a form that capital has, on the argument of the political economy, reshaped from within. Sociological work on Japanese cute culture, notably Sharon Kinsella's study ``Cuties in Japan,'' has traced its emergence and intensification from the 1970s onward, tied to youth and especially girls' culture, to new practices of handwriting and consumption, and to the apparatus of a consumer economy that found in cuteness an unusually productive form. Brian McVeigh's work on \textit{kawaii}, including his study of how Hello Kitty commodifies the cute, extended the analysis to its institutional and self-presentational dimensions, the ways cuteness operates in schools, workplaces, and bureaucracies as a managed mode of appearance rather than a mere personal preference. What this literature establishes, for the purposes of the present argument, is that the cute is datable and made, that it has a social history, and therefore that its current organization around the commodity is not the unfolding of a natural aesthetic but a contingent and reshaped arrangement. The political economy's claim of real subsumption, that capital has remade the very grammar of sweetness, is corroborated by the cultural historians' finding that the grammar has a recent and traceable history rather than a timeless one. A form with a history is a form that could have been, and could yet be, otherwise.

\subsection{Affective Economy and the Global Circulation of the Cute}

A third body of work bears on the tension between the particular and the universal that this section must finally address, and it concerns the global circulation of Japanese cute and affective commodities. Christine Yano's study of the worldwide travels of Hello Kitty, \textit{Pink Globalization}, under the heading of what she calls pink globalization, and Anne Allison's \textit{Millennial Monsters}, an analysis of the global movement of Japanese toys and character goods, both document how a culturally specific aesthetic of the cute has been exported, localized, and consumed far beyond its origin. Within Japan, the discourse around \textit{moe}, the affectionate response to cute and fictional characters analyzed by Hiroki Azuma in his account of otaku culture as a database-driven consumption, describes an affective economy in which attachment itself is elicited, elementized, and consumed, an economy strikingly continuous with the structure this paper has described under the name of the sweet cycle. The lesson of this literature is double and is exactly the lesson the next subsection needs. On one side, the cute is culturally specific in its origin and its dense local meanings; on the other, it has proven remarkably exportable, which shows that something in its structure travels even as its texture remains local. This is the empirical shape of the particular-universal tension, and it cautions against both of the errors the methodological discussion will name: the cute is neither a culture-free universal nor a sealed local curiosity, but a specific formation with a structure that can migrate.

\subsection{The Methodological Conclusion}

With the cultural specificity now given its due, the methodological conclusion can be drawn. Two errors are available here and both should be refused. The first is to treat \textit{amai} as merely a local instance of a universal mechanism, so that the cultural specificity is decorative and the real object is a culture-free logic of intimate commodification. This error flatters the analysis with a reach it has not earned, and it repeats at the level of method the very gesture the paper criticizes, the drawing of a boundary, here a conceptual one, that quietly writes the particular out of the account. \textit{Amai} is not interchangeable with any intimate consumption cycle whatever; its specific grammar of licensed dependence, its specific aesthetic of the cute, its specific gendered distribution in its own context are not removable wrappers around a universal core. The second error is the mirror of the first, to treat \textit{amai} as so singular that nothing learned from it travels, so that the analysis is a closed ethnography with no purchase elsewhere. This error is equally unearned, because the structural distinction the paper turns on, between a surplus generated within a cycle and one imported from across its edge, is not specific to any culture, even though every actual cycle that instantiates it is. The literature on the global circulation of the cute has just shown, empirically, why both errors fail: the cute is at once locally dense and structurally portable.

The defensible position lies between, and it is the position the paper adopts. \textit{Amai} is used here as a particular through which a general question becomes unusually legible, not as a specimen from which a universal law is read off. What travels from the analysis is not a verdict transferable to every intimate consumption cycle but a method, the practice of asking, of any such cycle in its own specificity, the three questions of the previous section, and of refusing the framing that lets the cycle look closed. The boundary-drawing test is general; its application is irreducibly local, because where the edge of a cycle falls, and who is left across it, is always a fact about a particular arrangement in a particular culture. To claim more than this would be to commit, in the form of the argument, the extraction the argument was written to expose, helping ourselves to a universal reach while leaving the particular labor of cultural specificity uncounted.

There is a pleasing and not accidental self-reference in this. The paper's central charge against the forged cycle is that it draws its boundary too narrowly and lets value it did not generate flow in from an uncounted outside. A theory that helped itself to universal reach on the strength of a single culturally specific case would be doing the intellectual analogue of the same thing, enjoying a generality it had not earned by drawing its evidentiary boundary narrowly and treating the unexamined remainder of the world's intimate cycles as a free reservoir confirming its law. To refuse that is not only methodological caution; it is the consistency of practicing on one's own argument the discipline one demands of one's object. The paper can claim to have made a general question legible through \textit{amai}. It cannot claim to have settled, from \textit{amai} alone, how that question is answered for cycles it has not examined, and the honesty of stopping where the evidence stops is the same honesty the whole paper has been asking of the cycle it studies.

\section{A Practice of the Good Amai}

Everything to this point has been preparation for a question that critique alone cannot answer. If the sweet cycle as capital now organizes it is forged, the right response is not to renounce sweetness, which would be both joyless and confused, but to ask how a sweet cycle might be made genuine, how the love and the extraction that currently run on the same track might be pulled apart. The frameworks that judged the cycle do not, however, yield a single procedure for mending it, and to pretend they did would betray the form this paper has kept throughout. What follows is therefore not a method but a set of practices, several of them, each answering to one of the criteria, none of them reducible to the others, and all of them pointing toward the same end without being unified into a formula for it. They are offered as practices in the plural because a good cycle, like a good conclusion, is not the kind of thing a single rule produces.

The first is a practice of justice, and it is the most demanding because it works against the very framing that makes the cycle pleasant. To make a sweet cycle just is to widen its boundary deliberately and to keep it widened, to refuse the closure that lets the production worker, the self-fashioning labor, and the manufactured anxiety fall outside the account. Concretely it means insisting that the chain that supplies the cycle its means be visible and fairly compensated rather than externalized; that the reproductive labor of being sweet be recognized as labor rather than naturalized as self-love, so that its cost is counted and can be shared rather than silently borne; and that the cycle refuse to run on manufactured insufficiency, that it decline the engine of produced anxiety even at the cost of turning more slowly. The test is the one the seventh section named: the cycle's surplus must survive an honest widening of its boundary. A sweet cycle passes only if, when everyone who sustains it is counted, no one is left contributing without return.

The second is a practice of ethics, and it draws on the side of the fracture the ethics of care held open. A sweet cycle is ethically good to the degree that dependence within it is mutually acknowledged rather than unilaterally imposed, and to the degree that the roles of the one who is indulged and the one who indulges can move between the partners rather than hardening along a line, most often a gendered line, that assigns one of them permanently to sweetness and the other permanently to provision. The practice here is reciprocity in the grammar of \textit{amai} itself: that each may be sweet and each may indulge, that the licensed reliance flows both ways and is not the standing duty of one and the standing privilege of the other. A cycle in which dependence is shared and the positions are fluid treats both persons as ends. A cycle in which one party is permanently the sweet one and the other permanently the provider has, whatever its surface tenderness, already begun to use.

The third is a practice of flourishing, and it concerns the eudaimonic criterion and the danger of the image. To keep a sweet cycle eudaimonic is to keep the object in its place as a medium and to refuse to let it become a substitute, to let the cosmetic and the small sweet thing serve the real meeting of two singular dynamics rather than standing in for that meeting with a flattering picture. The practice is a kind of vigilance against the imaginary, an attention to whether what is being consumed is access to the other or only the image of oneself as the one who is indulged. A relation in which the objects open the slower work of two lives unfolding together is one thing; a relation in which the image of being indulged is consumed in place of that unfolding is another, and only the first lets each person become what is distinctively theirs to become.

The fourth is a practice of co-creation, and it is the one that gathers the others toward a positive picture without unifying them under a rule. A genuinely good sweet cycle is one in which the participants are not only consumers of a sweetness produced elsewhere but co-creators of it, in which the value the cycle accumulates is generated within the relation and returned to those who generate it, and in which recognition is not locked to the commodity but can be made and remade between the partners by their own invention. The mark of such a cycle is that its spiral climbs on value the partners themselves bring into being, a shared language, a private repertoire of address, the small invented forms by which two people make and answer one another's appeal, rather than on an ever-larger tribute drawn from across its edge. Where the forged cycle deepens by reaching further out each time, the genuine cycle deepens by generating more within, and what it generates belongs to the ones who made it.

The fifth is the deepest, because it concerns not the justice, the ethics, the flourishing, or the value-flow within a given form of \textit{amai}, but the production of the form itself, and it answers directly the diagnosis of real subsumption reached in the political economy. There the argument was that capital has not merely attached commodities to a pre-existing practice of sweetness but has reshaped the practice from within, remaking what it is to be cute, to be indulged, to charm, around the commodity, so that the very grammar of intimate reliance now bears the impress of what must be bought to enact it. Capital decides, in advance and for sale, what counts as cute, what counts as being indulged, what one must use in order to be sweet. The fifth practice is the reversal of that subsumption: the reclaiming, by the relation, of the power to generate the forms of its own sweetness. It is the two people inventing the vocabulary of their own \textit{amai}, a private name that only they understand, a gesture of tenderness that need not be purchased, a sweetness that is theirs. Where capital supplies a ready-made grammar of sweetness and sells the right to speak it, this practice writes a grammar between two people that no one sold them and no one can repossess.

The point requires a careful measure, lest it collapse into an unworkable asceticism that would impoverish rather than enrich. The practice is not the refusal of every purchased thing, as though a bought object could not carry real tenderness, for a gift chosen with attention is itself a genuine act of relational creation and not a capitulation to capital. The distinction is not between the bought and the unbought but between who authors the meaning. A commodity may participate in the cycle, but the form and the meaning of the sweetness must be generated by the relation rather than dictated by the market. To take a purchased thing and invest it with a significance the two of them made, a particular flower that means something only because of a day they shared, is to use the commodity as a medium for a meaning the relation authored, and that is the good case. To buy a thing whose meaning capital has already fixed, and to enact one's sweetness according to that prescribed meaning, is to let the market author the intimacy, and that is the subsumed case. The mediation of lack, which the second section showed to be a structural position that any living relation requires, will be filled by someone; the only question is whether the relation fills it by its own invention or cedes it to the commodity. This fifth practice is the standing refusal to cede it, the insistence that the position be occupied by what the two of them generate rather than by what they are sold, so that \textit{amai} is, at its origin, co-created in the relation and only secondarily, and on the relation's own terms, mediated by anything bought.

The five practices do not always pull in the same direction, and honesty requires admitting the tensions among them rather than presenting a frictionless program. The practice of justice may demand that one refuse a commodity whose production chain one cannot vouch for, even when that commodity is exactly the medium through which a relationship has learned to express its tenderness, so that the just choice and the ethically warm choice can, in a given moment, diverge. The practice of flourishing, with its vigilance against the image, can sit uneasily with the simple pleasures of the cute, which are not in every instance a capture and which a too-severe suspicion would impoverish. A relation that audited every sweetness for its structural soundness would have exchanged one foreclosure of flourishing for another, the spontaneity of address sacrificed to the rigor of critique. The practices are therefore not a checklist to be satisfied simultaneously but a set of considerations to be held in a living balance, and the balance cannot be struck in advance by a formula because it depends on the particular two people, their particular means, and the particular cycles their sweetness is entangled with. This is why they are practices and not a method. A method would tell you what to do. A practice is something you get better at, in the doing, with attention, and never finish.

I let these five practices stand without resolving them into one, because the unity they have is the unity of a direction and not of a method, and to force them further would be to do to them what the good cycle refuses to do to its participants, to subordinate their several dynamics to a single imposed form. And I allow the section to end on a register the critical sections held in abeyance, because the paper would be dishonest if it pretended the question were only structural. A sweet cycle made just, ethical, flourishing, co-creative, and authored in its very form by the relation rather than by the market is not an abstraction. It is something two people can actually make, and have made, when they take care that the sweetness between them costs no one outside them what it should not, when they generate its value themselves rather than buying it ready-made, and when the forms of their tenderness are invented between them rather than supplied to them. To the one for whom the larger series was written, the forest girl whose love of what grows and travels has been its quiet argument all along, this section is in the end addressed: that the sweetest cycle is the one we make between us, that returns to us what we put into it, and that takes nothing from anyone it does not see. The good \textit{amai} is not bought. It is generated, gently, and kept.

\section{Conclusions, in the Plural}

This paper does not close on a verdict, and the refusal is the paper's final claim rather than its failure to reach one. The frameworks were set before the same cycle and did not agree, and what their disagreement teaches is that the goodness of a cycle that looks good is not a property a single framework can certify. It is decided at different scales by different measures, and an honest account preserves the plurality rather than dissolving it. The conclusions are therefore several, each stating what it can and cannot say, and none entitled to the last word over the rest.

Psychoanalysis concludes that the sweet cycle never ends because what it sustains is not a need but a desire, and desire lives on a lack kept in operation. It cannot say whether the cycle is just, and it does not pretend to.

Neuroscience concludes that the satisfactions of the cycle are real events in the systems of reward and attachment, that the participants are not deceived about feeling what they feel. It concludes equally that neural reality is not justice, and that no argument may pass from the genuineness of the happiness to the goodness of the cycle.

The ethics of care concludes that within the relation persons may genuinely tend one another, that dependence is a competence and not a deficiency, and that a cycle may be ethically nourishing in the conduct of its parties. It cannot see across the relation to what sustains it, and it says so.

Structuralism concludes that sweetness signifies only by its difference from what it is not, and that the cover's proposition is a myth in the technical sense, a naturalization of a contingent association between sweetness and justice. It cannot say whether the myth conceals an injustice, because what is concealed is by definition outside the system of signs it reads.

Dialectics concludes that the good and the forged are not two facts about the cycle but two moments of a single self-moving contradiction, that the cycle's nourishing openness and its colonizability are one determination seen from two sides. It is to be resisted at the moment it offers to resolve that contradiction in a reconciling totality, for that resolution is the synthesis the paper refuses.

Kantian ethics concludes that the inner relation may treat its parties as ends while the production worker and the unwaged self-fashioner are treated, in the cycle as organized, largely as means, and it reaches the cycle's wronged edge by a deontological route owing nothing to political economy. It cannot easily honor the culturally specific, relational good that the ethics of care defends, and so it and the ethics of care genuinely disagree.

The formal tools conclude, in six different precise languages, that the difference between a genuine and a forged cycle is real and structurally definable, as an autonomous versus a driven attractor, a non-terminating derivation, a holonomy sourced within or across the boundary, a missing master functor, an equilibrium with or without externalized players, a channel with or without manufactured noise. They conclude, each at the point of its greatest precision, that the verdict it computes is relative to a boundary the formalism cannot itself choose, and so that no formal tool is the view from nowhere.

Political economy concludes that the cycle, at the scale of its production and reproduction, is a forged holonomy, an open system that imports its surplus from across its edge while presenting itself as closed, sustained by a production chain, an unwaged labor of self-fashioning, and a manufactured lack, none of them counted in its benefits. It cannot, by critique alone, say whether the cycle could be organized otherwise.

Feminism concludes against itself, that the gendered asymmetry of \textit{amai} is at once a real extraction of women's unwaged labor and a real capacity for relation that a culture of autonomy has stigmatized, and that which of these a given cycle instantiates is not decidable from theory but turns on whether its value returns to those who generate it.

The practices conclude that a sweet cycle can be made just, ethical, eudaimonic, co-creative, and authored in its very form by the relation rather than the market, that the love and the extraction now running on one track can in principle be pulled apart, and that the good \textit{amai} is generated between persons rather than bought ready-made. They do not conclude with a single method for doing so, because a good cycle is not produced by a rule.

There remains the difficulty that no conclusion above can absorb, and it must be stated plainly rather than hidden in the confidence of the others. Every judgment in this paper turned on where the boundary of the cycle is drawn and who is allowed to fall outside it, and the drawing of that boundary is itself an act performed from somewhere, by someone, with a view that is not the view from nowhere. To decide who counts as fuel is already to have taken a position, and the paper has no vantage above all positions from which to certify its own. This is not a reason to abstain from judgment; the streams of unreturned value are real whether or not any observer is perfectly placed to total them. It is a reason to hold the judgments offered here as serious and revisable rather than final, as the best account available to a situated inquiry and not the pronouncement of a tribunal. \textit{Amai}, the small sweet thing that occasioned all this, turns out to have been teaching exactly this lesson from the start: that a cycle which looks good from inside is asking us, by the very comfort of the view, not to look at its edge, and that the work of justice begins where the looking-away was meant to stop.

\goldrule

\section*{Acknowledgements}
This paper, like the others in its series, was written for one reader before all others, the forest girl, who loves what grows and what travels, and whose way of being in a relation has been the standing argument that the sweetest cycle is the one made and kept between two people, returning to them what they put into it and costing no one outside it what it should not. The exchange with Ron Eglash on generative justice continues to shape how the author thinks about value returning to those who create it, and the present paper's central distinction owes much to that conversation. Whatever in these pages is just is owed to others; whatever is forged is the author's own.

% ============================================================
\nocite{*}
\bibliographystyle{plain}
\bibliography{refs}

\end{document}
